\documentclass[12pt]{article}

% === CÀI ĐẶT TRANG ===
\usepackage[margin=2cm]{geometry}
\usepackage[utf8]{inputenc}
\usepackage[T5]{fontenc}
\usepackage[vietnamese]{babel}

% === TOÁN HỌC ===
\usepackage{amsmath, amssymb, amsthm}

% === HÌNH VẼ ===
\usepackage{graphicx}
\usepackage{float}
\usepackage{tikz}
\usepackage{pgfplots}
\pgfplotsset{compat=1.18}

% === ĐỊNH DẠNG ===
\usepackage{enumitem}
\usepackage{xcolor}
\usepackage[most]{tcolorbox}
\usepackage{multicol}
\usepackage{booktabs}
\usepackage{array}
\usepackage{fancyhdr}
\usepackage{tocloft}
\usepackage[hidelinks]{hyperref}

% === LỆNH TỰ ĐỊNH NGHĨA ===
\newcommand{\otraloi}{\underline{\hspace{5cm}}}

% === TCOLORBOX STYLES ===
\tcbuselibrary{breakable, skins, fitting}

% Ví dụ
\newtcolorbox{vidu}[1][1]{
	colback=blue!5!white, colframe=blue!60!black,
	fonttitle=\bfseries, title={Ví dụ #1},
	breakable, enhanced
}

% Lời giải
\newtcolorbox{loigiai}{
	colback=gray!8!white, colframe=gray!50!black,
	fonttitle=\bfseries, title={Lời giải},
	breakable, enhanced
}

% Bài tập xanh – Bắt buộc
\newtcolorbox{btxanh}[1][]{
	colback=blue!6!white, colframe=blue!65!black,
	fonttitle=\bfseries\color{white}, coltitle=white,
	attach boxed title to top left={yshift=-2mm, xshift=4mm},
	boxed title style={colback=blue!65!black, rounded corners},breakable, enhanced, #1
}

% Bài tập vàng – Bổ sung
\newtcolorbox{btvang}[1][]{
	colback=yellow!12!white, colframe=orange!75!black,
	fonttitle=\bfseries, coltitle=orange!80!black,
	attach boxed title to top left={yshift=-2mm, xshift=4mm},
	boxed title style={colback=yellow!50!white, colframe=orange!75!black, rounded corners},breakable, enhanced, #1
}

% Bài tập đỏ – Gia cố
\newtcolorbox{btdo}[1][]{
	colback=red!6!white, colframe=red!65!black,
	fonttitle=\bfseries\color{white}, coltitle=white,
	attach boxed title to top left={yshift=-2mm, xshift=4mm},
	boxed title style={colback=red!65!black, rounded corners},breakable, enhanced, #1
}

% === HEADER/FOOTER ===
\pagestyle{fancy}
\fancyhf{}
\fancyhead[L]{\small \textbf{Giáo án: vector trong không gian}}
\fancyhead[R]{\small Lớp: 12}
\fancyfoot[C]{\thepage}
\renewcommand{\headrulewidth}{0.4pt}

% ================================================
\begin{document}
	
	\begin{center}
		{\Large\bfseries CHUYÊN ĐỀ II -- VECTOR TRONG KHÔNG GIAN}\\
		{\Large \textbf{BÀI 1.1: Vector trong không gian}}\\[4pt]
		{\large Môn: Toán \quad Lớp: 12}
	\end{center}
	\vspace{4mm}
	\hrule
	\vspace{4mm}
	
	% ===========================
	% PHẦN 1: MỤC TIÊU BÀI HỌC
	% ===========================
	\section*{Mục tiêu bài học}
	\addcontentsline{toc}{section}{Mục tiêu bài học}
	
	\textbf{1. Kiến thức:}
	\begin{itemize}
		\item Nắm được khái niệm vector, độ dài vector, vector cùng phương, vector bằng nhau trong không gian.
		\item Hiểu và vận dụng được các phép toán vector: tổng, hiệu, tích của một số với một vector trong không gian.
		\item Hiểu và vận dụng được tích vô hướng của hai vector, góc giữa hai vector trong không gian.
	\end{itemize}
	
	\noindent \textbf{2. Kĩ năng:}
	\begin{itemize}
		\item Xác định được góc giữa hai vector, tính được tích vô hướng.
		\item Biểu diễn, phân tích được một vector theo ba vector không đồng phẳng.
		\item Vận dụng các tính chất, quy tắc (quy tắc 3 điểm, quy tắc hình bình hành, quy tắc hình hộp) để chứng minh các đẳng thức vector.
	\end{itemize}
	
	\noindent \textbf{3. Thái độ / Phẩm chất:}
	\begin{itemize}
		\item Nghiêm túc, cẩn thận trong việc vẽ hình và tính toán.
		\item Phát triển tư duy không gian và khả năng logic.
	\end{itemize}
	
	\newpage
	
	% =====================
	% PHẦN 2: MỤC LỤC
	% =====================
	\setcounter{tocdepth}{2}
	\tableofcontents
	\newpage
	
	% =====================
	% PHẦN 3: LÝ THUYẾT
	% =====================
	\section{Lý thuyết}
	
	\subsection{Khái niệm vector trong không gian}
	
	\begin{itemize}
		\item \textbf{Định nghĩa:} vector trong không gian là một đoạn thẳng có hướng. Kí hiệu vector có điểm đầu $A$, điểm cuối $B$ là $\overrightarrow{AB}$. Nếu không cần chỉ rõ điểm đầu và điểm cuối, vector còn được kí hiệu là $\vec{a}, \vec{b}, \vec{x}, \vec{y}, \dots$
		\item \textbf{Giá của vector:} Đường thẳng đi qua điểm đầu và điểm cuối của một vector được gọi là giá của vector đó.
	\end{itemize}
	
	\begin{figure}[H]
		\centering
		\begin{tikzpicture}
			\draw[thick, color=gray] (0,0) -- (6,3) node[right, text=black] {$d$};
			\draw[-stealth, thick, blue] (1.5,0.75) -- (4.5,2.25) node[midway, above left, text=black] {$\vec{a}$};
		\end{tikzpicture}
		\caption{Đường thẳng $d$ là giá của vector $\vec{a}$}
	\end{figure}
	
	\begin{itemize}
		\item \textbf{Độ dài:} Là khoảng cách giữa điểm đầu và điểm cuối của vector đó. Kí hiệu độ dài của $\overrightarrow{AB}$ là $|\overrightarrow{AB}|$, của $\vec{a}$ là $|\vec{a}|$.
		\item \textbf{vector cùng phương:} Hai vector được gọi là cùng phương nếu chúng có giá song song hoặc trùng nhau. Nếu hai vector cùng phương thì chúng có thể cùng hướng hoặc ngược hướng.
		\item \textbf{vector bằng nhau:} Hai vector $\vec{a}$ và $\vec{b}$ được gọi là bằng nhau, kí hiệu $\vec{a} = \vec{b}$, nếu chúng có cùng độ dài và cùng hướng.
		\item \textbf{vector-không:} Là các vector có điểm đầu và điểm cuối trùng nhau (VD: $\overrightarrow{AA}, \overrightarrow{BB}, \dots$). vector-không có độ dài bằng $0$, cùng hướng (và cùng phương) với mọi vector. Kí hiệu chung là $\vec{0}$.
	\end{itemize}
	
	\begin{figure}[H]
		\centering
		\begin{tikzpicture}
			% Vẽ vector a
			\draw[-stealth, thick] (0,1) -- (2,2) node[midway, above left] {$\vec{a}$};
			% Vẽ vector b bằng a
			\draw[-stealth, thick] (3,1) -- (5,2) node[midway, above left] {$\vec{b}$};
			% Vẽ vector c cùng phương, ngược hướng
			\draw[-stealth, thick] (6,2.5) -- (4,1.5) node[midway, below right] {$\vec{c}$};
		\end{tikzpicture}
		\caption{Hai vector $\vec{a}$ và $\vec{b}$ bằng nhau; $\vec{a}$ và $\vec{c}$ cùng phương nhưng ngược hướng}
	\end{figure}
	
	\begin{vidu}[1]
		Cho hình hộp $ABCD.A'B'C'D'$. 
		\begin{center}
			\begin{tikzpicture}[scale=0.8]
				\coordinate (A) at (0,0);
				\coordinate (B) at (3,0);
				\coordinate (C) at (4,1.5);
				\coordinate (D) at (1,1.5);
				\coordinate (A') at (0,3);
				\coordinate (B') at (3,3);
				\coordinate (C') at (4,4.5);
				\coordinate (D') at (1,4.5);
				
				\draw[dashed] (A) -- (D) -- (C);
				\draw[dashed] (D) -- (D');
				\draw (A) -- (B) -- (C) -- (C') -- (B') -- (A') -- cycle;
				\draw (B) -- (B');
				\draw (A') -- (D') -- (C');
				
				\node[below left] at (A) {$A$};
				\node[below right] at (B) {$B$};
				\node[right] at (C) {$C$};
				\node[above left] at (D) {$D$};
				\node[left] at (A') {$A'$};
				\node[right] at (B') {$B'$};
				\node[above right] at (C') {$C'$};
				\node[above left] at (D') {$D'$};
				
				\draw[-stealth, thick, blue] (A) -- (B);
			\end{tikzpicture}
		\end{center}
		\begin{enumerate}[label=\alph*)]
			\item Hãy kể tên các vector khác $\vec{0}$ bằng với vector $\overrightarrow{AB}$.
			\item vector $\overrightarrow{AC}$ và $\overrightarrow{A'C'}$ có bằng nhau không? Vì sao?
		\end{enumerate}
	\end{vidu}
	
	\begin{loigiai}
		\begin{enumerate}[label=\alph*)]
			\item Các mặt của hình hộp là các hình bình hành. Do đó, ta có các đoạn thẳng song song và bằng nhau: $AB = DC = A'B' = D'C'$.
			
			Suy ra các vector cùng hướng và cùng độ dài với $\overrightarrow{AB}$ là: $\overrightarrow{DC}$, $\overrightarrow{A'B'}$, $\overrightarrow{D'C'}$.
			
			Vậy: $\overrightarrow{AB} = \overrightarrow{DC} = \overrightarrow{A'B'} = \overrightarrow{D'C'}$.
			
			\item Xét tứ giác $ACC'A'$ có $AA' \parallel CC'$ và $AA' = CC'$ (vì cùng song song và bằng $BB'$).
			
			Suy ra $ACC'A'$ là hình bình hành. Do đó đoạn $AC$ song song và bằng $A'C'$. 
			
			Mặt khác, $\overrightarrow{AC}$ và $\overrightarrow{A'C'}$ cùng hướng. Vậy $\overrightarrow{AC} = \overrightarrow{A'C'}$.
		\end{enumerate}
	\end{loigiai}
	
	\subsection{Tổng và hiệu của hai vector trong không gian}
	
	\textbf{1. Tổng của các vector}
	\begin{itemize}
		\item Phép cộng vector trong không gian được thực hiện hoàn toàn tương tự như trong mặt phẳng.
		\item Các tính chất: giao hoán ($\vec{a} + \vec{b} = \vec{b} + \vec{a}$), kết hợp $\big((\vec{a} + \vec{b}) + \vec{c} = \vec{a} + (\vec{b} + \vec{c})\big)$ và cộng với vector-không ($\vec{a} + \vec{0} = \vec{a}$).
		\item Để thực hiện phép cộng, ta thường sử dụng các quy tắc sau:
	\end{itemize}
	
	\textbf{a) Quy tắc 3 điểm:} Với ba điểm $A, B, C$ bất kì, ta luôn có: 
	$$ \overrightarrow{AB} + \overrightarrow{BC} = \overrightarrow{AC} $$
	
	\begin{figure}[H]
		\centering
		\begin{tikzpicture}
			\coordinate (A) at (0,0);
			\coordinate (B) at (2,1.5);
			\coordinate (C) at (5,0.5);
			
			\draw[-stealth, thick, blue] (A) -- (B);
			\draw[-stealth, thick, blue] (B) -- (C);
			\draw[-stealth, thick, red] (A) -- (C);
			
			\node[left] at (A) {$A$};
			\node[above] at (B) {$B$};
			\node[right] at (C) {$C$};
		\end{tikzpicture}
		\caption{Quy tắc 3 điểm}
	\end{figure}
	
	\textbf{b) Quy tắc hình bình hành:} Nếu $ABCD$ là hình bình hành thì:
	$$ \overrightarrow{AB} + \overrightarrow{AD} = \overrightarrow{AC} $$
	
	\begin{figure}[H]
		\centering
		\begin{tikzpicture}
			\coordinate (A) at (0,0);
			\coordinate (B) at (4,0);
			\coordinate (C) at (5,2);
			\coordinate (D) at (1,2);
			
			\draw[dashed] (A) -- (D) -- (C);
			\draw (A) -- (B) -- (C);
			
			\draw[-stealth, thick, blue] (A) -- (B);
			\draw[-stealth, thick, blue] (A) -- (D);
			\draw[-stealth, thick, red] (A) -- (C);
			
			\node[below left] at (A) {$A$};
			\node[below right] at (B) {$B$};
			\node[above right] at (C) {$C$};
			\node[above left] at (D) {$D$};
		\end{tikzpicture}
		\caption{Quy tắc hình bình hành}
	\end{figure}
	
	\textbf{c) Quy tắc hình hộp:} Cho hình hộp $ABCD.A'B'C'D'$. Ta có:
	$$ \overrightarrow{AB} + \overrightarrow{AD} + \overrightarrow{AA'} = \overrightarrow{AC'} $$
	
	\begin{figure}[H]
		\centering
		\begin{tikzpicture}[scale=1.2, line cap=round, line join=round]
			% Định nghĩa tọa độ các đỉnh để tạo góc nhìn tách biệt các vector
			\coordinate (A)  at (0,0);
			\coordinate (B)  at (3,-0.5);
			\coordinate (D)  at (2.5,1);
			\coordinate (A') at (0,3);
			
			\coordinate (C)  at (5.5,0.5);
			\coordinate (B') at (3,2.5);
			\coordinate (D') at (2.5,4);
			\coordinate (C') at (5.5,3.5);
			
			% 1. Vẽ các cạnh khuất ở phía sau
			\draw[dashed] (D) -- (C);
			\draw[dashed] (D) -- (D');
			
			% 2. Vẽ vector thành phần nằm mặt sau (AD) - nét đứt
			\draw[-stealth, thick, blue, dashed] (A) -- (D);
			
			% 3. Vẽ vector tổng (AC') xuyên qua lòng hình hộp
			\draw[-stealth, thick, red] (A) -- (C');
			
			% 4. Vẽ các cạnh liền của hình hộp (vẽ sau để đè lên nét xiên AC' tạo ảo giác 3D)
			\draw (B) -- (C) -- (C') -- (D') -- (A') -- (B') -- cycle;
			\draw (B) -- (B') -- (C');
			
			% 5. Vẽ 2 vector thành phần nằm mặt trước (AB, AA')
			\draw[-stealth, thick, blue] (A) -- (B);
			\draw[-stealth, thick, blue] (A) -- (A');
			
			% 6. Nhãn các đỉnh
			\node[below left] at (A) {$A$};
			\node[below] at (B) {$B$};
			\node[right] at (C) {$C$};
			\node[above left] at (D) {$D$};
			\node[left] at (A') {$A'$};
			\node[above left] at (B') {$B'$};
			\node[above right] at (C') {$C'$};
			\node[above] at (D') {$D'$};
		\end{tikzpicture}
		\caption{Quy tắc hình hộp: $\overrightarrow{AB} + \overrightarrow{AD} + \overrightarrow{AA'} = \overrightarrow{AC'}$}
	\end{figure}
	
	\vspace{0.5cm}
	\textbf{2. Hiệu của hai vector}
	\begin{itemize}
		\item \textbf{vector đối:} vector có cùng độ dài và ngược hướng với vector $\vec{a}$ được gọi là vector đối của $\vec{a}$, kí hiệu là $-\vec{a}$. Mọi vector $\overrightarrow{AB}$ đều có vector đối là $\overrightarrow{BA}$, tức là $\overrightarrow{AB} = -\overrightarrow{BA}$.
		\item \textbf{Phép trừ vector:} $\vec{a} - \vec{b} = \vec{a} + (-\vec{b})$.
		\item \textbf{Quy tắc hiệu (quy tắc 3 điểm cho phép trừ):} Với ba điểm $O, A, B$ bất kì, ta có:
		$$ \overrightarrow{OB} - \overrightarrow{OA} = \overrightarrow{AB} $$
	\end{itemize}
	
	\begin{figure}[H]
		\centering
		\begin{tikzpicture}
			\coordinate (O) at (0,0);
			\coordinate (A) at (2,2);
			\coordinate (B) at (4,0.5);
			
			\draw[-stealth, thick, blue] (O) -- (A);
			\draw[-stealth, thick, blue] (O) -- (B);
			\draw[-stealth, thick, red] (A) -- (B);
			
			\node[left] at (O) {$O$};
			\node[above] at (A) {$A$};
			\node[right] at (B) {$B$};
		\end{tikzpicture}
		\caption{Quy tắc hiệu}
	\end{figure}
	
	\begin{vidu}[2]
		Cho hình hộp $ABCD.A'B'C'D'$. Chứng minh các đẳng thức vector sau:
		\begin{enumerate}[label=\alph*)]
			\item $\overrightarrow{AB} + \overrightarrow{B'C'} + \overrightarrow{DD'} = \overrightarrow{AC'}$
			\item $\overrightarrow{AC'} - \overrightarrow{AA'} - \overrightarrow{AD} = \overrightarrow{AB}$
		\end{enumerate}
	\end{vidu}
	
	\begin{loigiai}
		\begin{enumerate}[label=\alph*)]
			\item Vì $ABCD.A'B'C'D'$ là hình hộp nên các mặt đối diện là các hình bình hành bằng nhau. Do đó ta có:
			$\overrightarrow{B'C'} = \overrightarrow{BC}$ và $\overrightarrow{DD'} = \overrightarrow{CC'}$.
			
			Biến đổi vế trái (VT) và áp dụng quy tắc 3 điểm liên tiếp:
			$$ \text{VT} = \overrightarrow{AB} + \overrightarrow{B'C'} + \overrightarrow{DD'} = \overrightarrow{AB} + \overrightarrow{BC} + \overrightarrow{CC'} $$
			$$ \text{VT} = (\overrightarrow{AB} + \overrightarrow{BC}) + \overrightarrow{CC'} = \overrightarrow{AC} + \overrightarrow{CC'} = \overrightarrow{AC'} = \text{VP} \text{ (điều phải chứng minh)} $$
			
			\item Áp dụng quy tắc hình hộp ta có: $\overrightarrow{AB} + \overrightarrow{AD} + \overrightarrow{AA'} = \overrightarrow{AC'}$.
			
			Từ đẳng thức trên, ta chuyển vế các vector $\overrightarrow{AA'}$ và $\overrightarrow{AD}$ (đổi dấu thành phép trừ):
			$$ \overrightarrow{AC'} - \overrightarrow{AA'} - \overrightarrow{AD} = \overrightarrow{AB} \text{ (điều phải chứng minh)} $$
		\end{enumerate}
	\end{loigiai}
	
	
	
	\subsection{Tích của một số với một vector trong không gian}
	
	\subsubsection{Định nghĩa}
	\begin{itemize}
		\item Trong không gian, tích của một số thực $k \neq 0$ với một vector $\vec{a} \neq \vec{0}$ là một vector, kí hiệu là $k\vec{a}$, được xác định như sau:
		\begin{itemize}
			\item Cùng hướng với $\vec{a}$ nếu $k > 0$; ngược hướng với $\vec{a}$ nếu $k < 0$.
			\item Có độ dài bằng $|k| \cdot |\vec{a}|$.
		\end{itemize}
		\item \textbf{Quy ước:} $0\vec{a} = \vec{0}$ và $k\vec{0} = \vec{0}$.
		\item \textbf{Điều kiện cùng phương:} Hai vector $\vec{a}$ và $\vec{b}$ ($\vec{b} \neq \vec{0}$) cùng phương khi và chỉ khi tồn tại số thực $k$ sao cho $\vec{a} = k\vec{b}$.
	\end{itemize}
	
	\begin{figure}[H]
		\centering
		\begin{tikzpicture}
			% vector gốc a
			\draw[-stealth, thick, blue] (0,0) -- (2,1) node[midway, above left] {$\vec{a}$};
			
			% vector k*a với k > 0
			\draw[-stealth, thick, red] (3,0) -- (7,2) node[midway, above left] {$k\vec{a} \ (k > 0)$};
			
			% vector k*a với k < 0
			\draw[-stealth, thick, orange] (7,-1) -- (4,-2.5) node[midway, below right] {$k\vec{a} \ (k < 0)$};
		\end{tikzpicture}
		\caption{Minh họa tích của một số với một vector}
	\end{figure}
	
	\subsubsection{Các tính chất}
	Với $h, k$ là các số thực và $\vec{a}, \vec{b}$ là các vector bất kì, ta có:
	\begin{itemize}
		\item Tính chất kết hợp: $h(k\vec{a}) = (hk)\vec{a}$.
		\item Tính chất phân phối: $(h+k)\vec{a} = h\vec{a} + k\vec{a}$ và $k(\vec{a} + \vec{b}) = k\vec{a} + k\vec{b}$.
		\item $1\vec{a} = \vec{a}$ và $(-1)\vec{a} = -\vec{a}$.
	\end{itemize}
	
	\subsubsection{Quy tắc trung điểm và quy tắc trọng tâm tam giác}
	
	\textbf{a) Quy tắc trung điểm:}
	Cho đoạn thẳng $AB$ có $I$ là trung điểm.
	\begin{itemize}
		\item \textbf{Hệ thức gốc:}
		$$ \overrightarrow{IA} + \overrightarrow{IB} = \vec{0} $$
		\item \textbf{Hệ thức mở rộng:} Với mọi điểm $M$ bất kì trong không gian, ta luôn có:
		$$ \overrightarrow{MA} + \overrightarrow{MB} = 2\overrightarrow{MI} $$
	\end{itemize}
	
	% === HÌNH 1: QUY TẮC TRUNG ĐIỂM (ĐỊNH LÝ GỐC) ===
	\begin{figure}[H]
		\centering
		\begin{tikzpicture}[scale=1.2, line cap=round, line join=round]
			% Tọa độ các điểm
			\coordinate (I) at (0,0);
			\coordinate (A) at (-2.5,0);
			\coordinate (B) at (2.5,0);
			
			% Vẽ đoạn thẳng AB
			\draw[thick, gray] (A) -- (B);
			
			% Vẽ các vector IA và IB xuất phát từ I
			\draw[-stealth, very thick, blue] (I) -- (A) node[midway, above] {$\overrightarrow{IA}$};
			\draw[-stealth, very thick, blue] (I) -- (B) node[midway, above] {$\overrightarrow{IB}$};
			
			% Vẽ ký hiệu trung điểm (vạch chia bằng nhau)
			\draw[thick] (-1.25,-0.15) -- (-1.25,0.15);
			\draw[thick] (1.25,-0.15) -- (1.25,0.15);
			
			% Điểm và nhãn
			\fill (I) circle (2pt) node[below] {$I$};
			\fill (A) circle (2pt) node[below] {$A$};
			\fill (B) circle (2pt) node[below] {$B$};
		\end{tikzpicture}
		\caption{Hệ thức gốc trung điểm: $\overrightarrow{IA} + \overrightarrow{IB} = \vec{0}$}
	\end{figure}
	
	% === HÌNH 2: QUY TẮC TRUNG ĐIỂM (HỆ THỨC MỞ RỘNG) ===
	\begin{figure}[H]
		\centering
		\begin{tikzpicture}[scale=1.2, line cap=round, line join=round]
			% Tọa độ
			\coordinate (I) at (2,0);
			\coordinate (A) at (0,0);
			\coordinate (B) at (4,0);
			\coordinate (M) at (1,2);
			\coordinate (M') at (3,-2); % M' sao cho I là trung điểm MM' (MM' = 2MI)
			
			% Vẽ đoạn thẳng đáy AB
			\draw[thick, gray] (A) -- (B);
			
			% Vẽ hình bình hành MAM'B bằng nét đứt để minh họa phép cộng vector
			\draw[dashed, gray] (A) -- (M') -- (B);
			
			% Vẽ các vector MA, MB (Xanh)
			\draw[-stealth, thick, blue] (M) -- (A) node[midway, left] {$\overrightarrow{MA}$};
			\draw[-stealth, thick, blue] (M) -- (B) node[midway, above right] {$\overrightarrow{MB}$};
			
			% Vẽ vector MI (Đỏ) và phần kéo dài IM' để minh họa 2MI
			\draw[-stealth, very thick, red] (M) -- (I) node[midway, left] {$\overrightarrow{MI}$};
			\draw[dashed, thick, red] (I) -- (M') node[midway, left] {};
			
			% Đánh dấu trung điểm I
			\draw[thick] (1,-0.1) -- (1,0.1);
			\draw[thick] (3,-0.1) -- (3,0.1);
			
			% Nhãn các điểm
			\fill (M) circle (2pt) node[above] {$M$};
			\fill (A) circle (2pt) node[left] {$A$};
			\fill (B) circle (2pt) node[right] {$B$};
			\fill (I) circle (2pt) node[above right] {$I$};
			\fill (M') circle (1.5pt) node[below] {$M'$};
		\end{tikzpicture}
		\caption{Hệ thức mở rộng trung điểm: $\overrightarrow{MA} + \overrightarrow{MB} = 2\overrightarrow{MI}$}
	\end{figure}
	
	
	\noindent \textbf{b) Quy tắc trọng tâm tam giác:}
	Cho tam giác $ABC$ có $G$ là trọng tâm.
	\begin{itemize}
		\item \textbf{Hệ thức gốc:}
		$$ \overrightarrow{GA} + \overrightarrow{GB} + \overrightarrow{GC} = \vec{0} $$
		\item \textbf{Hệ thức mở rộng:} Với mọi điểm $M$ bất kì trong không gian, ta luôn có:
		$$ \overrightarrow{MA} + \overrightarrow{MB} + \overrightarrow{MC} = 3\overrightarrow{MG} $$
	\end{itemize}
	
	\subsubsection{Sự đồng phẳng của ba vector}
	\begin{itemize}
		\item \textbf{Khái niệm:} Ba vector trong không gian được gọi là đồng phẳng nếu các giá của chúng cùng song song với một mặt phẳng.
		\item \textbf{Điều kiện đồng phẳng:} Cho ba vector $\vec{a}, \vec{b}, \vec{c}$ trong đó $\vec{a}, \vec{b}$ không cùng phương. Ba vector $\vec{a}, \vec{b}, \vec{c}$ đồng phẳng khi và chỉ khi tồn tại duy nhất cặp số thực $(m, n)$ sao cho:
		$$ \vec{c} = m\vec{a} + n\vec{b} $$
	\end{itemize}
	
	\begin{vidu}[3]
		Cho tứ diện $ABCD$. Gọi $G$ là trọng tâm của tứ diện (nghĩa là $G$ là trung điểm của đoạn thẳng nối trung điểm hai cạnh đối diện). 
		\begin{enumerate}[label=\alph*)]
			\item Chứng minh rằng: $\overrightarrow{GA} + \overrightarrow{GB} + \overrightarrow{GC} + \overrightarrow{GD} = \vec{0}$.
			\item Với điểm $M$ bất kì trong không gian, chứng minh: $\overrightarrow{MA} + \overrightarrow{MB} + \overrightarrow{MC} + \overrightarrow{MD} = 4\overrightarrow{MG}$.
		\end{enumerate}
	\end{vidu}
	
	\begin{loigiai}
		\begin{enumerate}[label=\alph*)]
			\item Gọi $I$ và $J$ lần lượt là trung điểm của các cạnh $AB$ và $CD$. Theo định nghĩa, $G$ là trung điểm của $IJ$.
			
			Áp dụng tính chất trung điểm, ta có:
			$$ \overrightarrow{GA} + \overrightarrow{GB} = 2\overrightarrow{GI} $$
			$$ \overrightarrow{GC} + \overrightarrow{GD} = 2\overrightarrow{GJ} $$
			Cộng vế theo vế, ta được:
			$$ \overrightarrow{GA} + \overrightarrow{GB} + \overrightarrow{GC} + \overrightarrow{GD} = 2(\overrightarrow{GI} + \overrightarrow{GJ}) $$
			Vì $G$ là trung điểm của $IJ$ nên $\overrightarrow{GI} + \overrightarrow{GJ} = \vec{0}$.
			
			Do đó: $\overrightarrow{GA} + \overrightarrow{GB} + \overrightarrow{GC} + \overrightarrow{GD} = \vec{0}$ (điều phải chứng minh).
			
			\item Theo quy tắc 3 điểm (chèn điểm $G$), ta có:
			$$ \overrightarrow{MA} = \overrightarrow{MG} + \overrightarrow{GA} $$
			$$ \overrightarrow{MB} = \overrightarrow{MG} + \overrightarrow{GB} $$
			$$ \overrightarrow{MC} = \overrightarrow{MG} + \overrightarrow{GC} $$
			$$ \overrightarrow{MD} = \overrightarrow{MG} + \overrightarrow{GD} $$
			Cộng vế theo vế 4 đẳng thức trên, ta được:
			$$ \overrightarrow{MA} + \overrightarrow{MB} + \overrightarrow{MC} + \overrightarrow{MD} = 4\overrightarrow{MG} + (\overrightarrow{GA} + \overrightarrow{GB} + \overrightarrow{GC} + \overrightarrow{GD}) $$
			Áp dụng kết quả của câu a, biểu thức trong ngoặc bằng $\vec{0}$.
			
			Vậy: $\overrightarrow{MA} + \overrightarrow{MB} + \overrightarrow{MC} + \overrightarrow{MD} = 4\overrightarrow{MG}$ (điều phải chứng minh).
		\end{enumerate}
	\end{loigiai}
	
		% === HÌNH 3: QUY TẮC TRỌNG TÂM TAM GIÁC (ĐỊNH LÝ GỐC) ===
	\begin{figure}[H]
		\centering
		\begin{tikzpicture}[scale=1.2, line cap=round, line join=round]
			% Tọa độ các đỉnh tam giác và trọng tâm G
			\coordinate (A) at (1.5,2.5);
			\coordinate (B) at (0,0);
			\coordinate (C) at (4,0.5);
			\coordinate (G) at (1.833,1); % G = ( (1.5+0+4)/3 , (2.5+0+0.5)/3 )
			
			% Vẽ khung tam giác ABC
			\draw[thick, gray] (A) -- (B) -- (C) -- cycle;
			
			% Vẽ các vector xuất phát từ G đến các đỉnh (Xanh)
			\draw[-stealth, very thick, blue] (G) -- (A) node[midway, left] {$\overrightarrow{GA}$};
			\draw[-stealth, very thick, blue] (G) -- (B) node[midway, below] {$\overrightarrow{GB}$};
			\draw[-stealth, very thick, blue] (G) -- (C) node[midway, above] {$\overrightarrow{GC}$};
			
			% Nhãn điểm
			\fill (G) circle (2pt) node[below right] {$G$};
			\fill (A) circle (2pt) node[above] {$A$};
			\fill (B) circle (2pt) node[left] {$B$};
			\fill (C) circle (2pt) node[right] {$C$};
		\end{tikzpicture}
		\caption{Hệ thức gốc trọng tâm tam giác: $\overrightarrow{GA} + \overrightarrow{GB} + \overrightarrow{GC} = \vec{0}$}
	\end{figure}
	
	\begin{figure}[H]
		\centering
		\begin{tikzpicture}[scale=1.3, line cap=round, line join=round]
			% --- ĐỊNH NGHĨA TỌA ĐỘ (Đã tối ưu hóa góc lệch để không trùng cạnh) ---
			\coordinate (M) at (-1.5, 4.5);
			\coordinate (A) at (2, 4);
			\coordinate (B) at (-0.5, 0.5);
			\coordinate (C) at (5, 1);
			
			% Tính toán tọa độ trọng tâm G chính xác
			\coordinate (G) at (2.167, 1.833); 
			
			% Điểm trung điểm BC để vẽ đường trung tuyến mờ
			\coordinate (I) at (2.25, 0.75);
			
			% --- VẼ HÌNH NỀN (Tam giác ABC và các đường dóng mờ) ---
			\draw[thick, gray!50, dashed] (A) -- (B) -- (C) -- cycle;
			\draw[gray!30, dotted] (A) -- (I); % Đường trung tuyến AI
			
			% --- VẼ CÁC vector THÀNH PHẦN (Màu xanh dương, tách biệt rõ ràng) ---
			\draw[-stealth, very thick, blue] (M) -- (A) 
			node[midway, above, yshift=2pt] {$\overrightarrow{MA}$};
			\draw[-stealth, very thick, blue] (M) -- (B) 
			node[midway, left, xshift=-2pt] {$\overrightarrow{MB}$};
			\draw[-stealth, very thick, blue] (M) -- (C) 
			node[midway, below, xshift=15pt, yshift=-2pt] {$\overrightarrow{MC}$};
			
			% --- VẼ vector TỔNG (Màu đỏ, nổi bật ở giữa) ---
			\draw[-stealth, very thick, red] (M) -- (G) 
			node[midway, left, xshift=-4pt, yshift=-8pt] {$\overrightarrow{MG}$};
			
			% --- ĐIỂM VÀ NHÃN ĐỈNH (Đặt lệch hướng để không đè lên nét vẽ) ---
			\fill (M) circle (2.5pt) node[above left] {$M$};
			\fill (A) circle (2pt) node[above right] {$A$};
			\fill (B) circle (2pt) node[left] {$B$};
			\fill (C) circle (2pt) node[right] {$C$};
			\fill (G) circle (2.5pt) node[below right] {$G$};
		\end{tikzpicture}
		\caption{Hệ thức mở rộng trọng tâm tam giác: $\overrightarrow{MA} + \overrightarrow{MB} + \overrightarrow{MC} = 3\overrightarrow{MG}$}
	\end{figure}
	
	
	\subsection{Tích vô hướng của hai vector trong không gian}
	
	\textbf{1. Góc giữa hai vector trong không gian}
	\begin{itemize}
		\item \textbf{Định nghĩa:} Trong không gian, cho hai vector $\vec{a}$ và $\vec{b}$ khác $\vec{0}$. Lấy một điểm $O$ bất kì, vẽ các vector $\overrightarrow{OA} = \vec{a}$ và $\overrightarrow{OB} = \vec{b}$. Góc $\widehat{AOB}$ với $0^\circ \le \widehat{AOB} \le 180^\circ$ được gọi là góc giữa hai vector $\vec{a}$ và $\vec{b}$, kí hiệu là $(\vec{a}, \vec{b})$.
		\item \textbf{Quy ước:} Nếu ít nhất một trong hai vector $\vec{a}$ hoặc $\vec{b}$ bằng $\vec{0}$ thì góc giữa chúng có thể nhận một giá trị tùy ý từ $0^\circ$ đến $180^\circ$.
		\item \textbf{vector vuông góc:} Khi $(\vec{a}, \vec{b}) = 90^\circ$, ta nói hai vector $\vec{a}$ và $\vec{b}$ vuông góc với nhau, kí hiệu là $\vec{a} \perp \vec{b}$.
	\end{itemize}
	
	\begin{figure}[H]
		\centering
		\begin{tikzpicture}[scale=1.2, line cap=round, line join=round]
			% Khởi tạo điểm O và vẽ hai tia OA, OB
			\coordinate (O) at (0,0);
			\coordinate (A) at (2.5,0.8);
			\coordinate (B) at (1.5,-1.2);
			
			% Vẽ vector OA, OB làm đại diện
			\draw[-stealth, thick, blue] (O) -- (A) node[above right] {$A$} node[midway, above left] {$\vec{a}$};
			\draw[-stealth, thick, blue] (O) -- (B) node[below right] {$B$} node[midway, below left] {$\vec{b}$};
			\fill (O) circle (2pt) node[left] {$O$};
			
			% Vẽ cung biểu diễn góc alpha giữa hai vector
			\draw (0.5, 0.16) arc (17.7:-38.6:0.5);
			\node at (0.8, -0.1) {$\alpha$};
			
			% Vẽ hai vector tự do ban đầu ngoài không gian để minh họa phép tịnh tiến
			\coordinate (U1) at (-2.5, 1);
			\coordinate (U2) at (-0.5, 1.8);
			\draw[-stealth, thick, gray] (U1) -- (U2) node[midway, above] {$\vec{a}$};
			
			\coordinate (V1) at (-3, -0.5);
			\coordinate (V2) at (-1.8, -2.1);
			\draw[-stealth, thick, gray] (V1) -- (V2) node[midway, left] {$\vec{b}$};
			
			% Các mũi tên đứt thể hiện phép tịnh tiến về O
			\draw[dashed, ->, gray!60] (U1) .. controls (-1.2,0.9) .. (O);
			\draw[dashed, ->, gray!60] (V1) .. controls (-1.8,-0.1) .. (O);
		\end{tikzpicture}
		\caption{Góc giữa hai vector $\vec{a}$ và $\vec{b}$ trong không gian}
	\end{figure}
	
	\textbf{2. Tích vô hướng của hai vector}
	\begin{itemize}
		\item \textbf{Định nghĩa:} Trong không gian, tích vô hướng của hai vector $\vec{a}$ và $\vec{b}$ khác $\vec{0}$ là một số (kí hiệu là $\vec{a} \cdot \vec{b}$) được xác định bởi công thức:
		$$ \vec{a} \cdot \vec{b} = |\vec{a}| \cdot |\vec{b}| \cdot \cos(\vec{a}, \vec{b}) $$
		\item \textbf{Quy ước:} Nếu $\vec{a} = \vec{0}$ hoặc $\vec{b} = \vec{0}$ thì $\vec{a} \cdot \vec{b} = 0$.
		\item \textbf{Các hệ quả quan trọng:}
		\begin{itemize}
			\item Điều kiện vuông góc: $\vec{a} \perp \vec{b} \Leftrightarrow \vec{a} \cdot \vec{b} = 0$.
			\item Bình phương vô hướng: $\vec{a}^2 = |\vec{a}|^2$ hay $|\vec{a}| = \sqrt{\vec{a}^2}$.
			\item Công thức tính góc: Nếu $\vec{a}, \vec{b}$ khác $\vec{0}$ thì ta có:
			$$ \cos(\vec{a}, \vec{b}) = \frac{\vec{a} \cdot \vec{b}}{|\vec{a}| \cdot |\vec{b}|} $$
		\end{itemize}
		\item \textbf{Tính chất:} Tương tự như trong mặt phẳng, tích vô hướng trong không gian cũng có các tính chất: giao hoán ($\vec{a}\cdot\vec{b} = \vec{b}\cdot\vec{a}$), phân phối ($\vec{a}\cdot(\vec{b} + \vec{c}) = \vec{a}\cdot\vec{b} + \vec{a}\cdot\vec{c}$), kết hợp số thực ($(k\vec{a})\cdot\vec{b} = k(\vec{a}\cdot\vec{b})$).
	\end{itemize}
	
	\begin{vidu}[4]
		Cho hình lập phương $ABCD.A'B'C'D'$ có cạnh bằng $a$.
		\begin{center}
			\begin{tikzpicture}[scale=1, line cap=round, line join=round]
				% Định nghĩa tọa độ trị số tuyệt đối trực tiếp (tránh dùng thư viện calc gây lỗi biên dịch)
				\coordinate (A) at (0,0);
				\coordinate (B) at (3,-0.5);
				\coordinate (D) at (2,1.2);
				\coordinate (A') at (0,3);
				
				\coordinate (C) at (5,0.7);
				\coordinate (B') at (3,2.5);
				\coordinate (D') at (2,4.2);
				\coordinate (C') at (5,3.7);
				
				% Vẽ các đường đứt nét phía sau
				\draw[dashed] (A) -- (D) -- (C);
				\draw[dashed] (D) -- (D');
				\draw[dashed, blue, thick] (A) -- (D);
				\draw[dashed, red, thick] (A) -- (C);
				
				% Vẽ khung ngoài hình lập phương
				\draw (A) -- (B) -- (C) -- (C') -- (B') -- (A') -- cycle;
				\draw (B) -- (B');
				\draw (A') -- (D') -- (C');
				\draw (B') -- (C');
				
				% vector AB nổi bật màu xanh
				\draw[blue, thick] (A) -- (B);
				
				% Đặt nhãn đỉnh
				\node[below left] at (A) {$A$};
				\node[below] at (B) {$B$};
				\node[right] at (C) {$C$};
				\node[above left] at (D) {$D$};
				\node[left] at (A') {$A'$};
				\node[right] at (B') {$B'$};
				\node[above right] at (C') {$C'$};
				\node[above] at (D') {$D'$};
			\end{tikzpicture}
		\end{center}
		\begin{enumerate}[label=\alph*)]
			\item Tính tích vô hướng của hai vector $\overrightarrow{AB}$ và $\overrightarrow{AD}$.
			\item Tính tích vô hướng của hai vector $\overrightarrow{AB}$ và $\overrightarrow{A'C'}$.
		\end{enumerate}
	\end{vidu}
	
	\begin{loigiai}
		\begin{enumerate}[label=\alph*)]
			\item Vì $ABCD.A'B'C'D'$ là hình lập phương nên mặt đáy $ABCD$ là hình vuông. Do đó, hai cạnh $AB$ và $AD$ vuông góc với nhau tại $A$, tức là $\overrightarrow{AB} \perp \overrightarrow{AD}$.
			
			Áp dụng điều kiện vuông góc của tích vô hướng:
			$$ \overrightarrow{AB} \cdot \overrightarrow{AD} = 0 $$
			
			\item Ta có tứ giác $ACC'A'$ là hình chữ nhật nên $\overrightarrow{A'C'} = \overrightarrow{AC}$.
			
			Do đó, tích vô hướng của hai vector $\overrightarrow{AB}$ và $\overrightarrow{A'C'}$ được đưa về tính trên mặt đáy $ABCD$:
			$$ \overrightarrow{AB} \cdot \overrightarrow{A'C'} = \overrightarrow{AB} \cdot \overrightarrow{AC} = |\overrightarrow{AB}| \cdot |\overrightarrow{AC}| \cdot \cos(\overrightarrow{AB}, \overrightarrow{AC}) $$
			Trong hình vuông $ABCD$ cạnh $a$:
			\begin{itemize}
				\item Độ dài cạnh: $AB = a$.
				\item Độ dài đường chéo: $AC = a\sqrt{2}$.
				\item Góc ở đỉnh hình vuông: $(\overrightarrow{AB}, \overrightarrow{AC}) = \widehat{BAC} = 45^\circ$.
			\end{itemize}
			Thay các giá trị vào công thức tích vô hướng:
			$$ \overrightarrow{AB} \cdot \overrightarrow{A'C'} = a \cdot a\sqrt{2} \cdot \cos(45^\circ) = a^2\sqrt{2} \cdot \frac{\sqrt{2}}{2} = a^2 $$
		\end{enumerate}
	\end{loigiai}
	
	\newpage
	
	% =====================
	% PHẦN 4: BÀI TẬP
	% =====================
	\section{Bài tập tự luận}

	
	\subsection{Dạng 1: Xác định vector và các tính chất của vector}	
	
	\begin{btxanh}[]
		Cho hình hộp $ABCD.A'B'C'D'$. Sử dụng các đỉnh của hình hộp làm điểm đầu và điểm cuối của các vectơ.
		\begin{enumerate}[label=\alph*)]
			\item Hãy kể tên các vectơ bằng nhau lần lượt bằng các vectơ $\overrightarrow{AB}, \overrightarrow{AC}, \overrightarrow{AD}, \overrightarrow{AA'}$.
			\item Hãy kể tên các vectơ luôn có độ dài bằng nhau và bằng độ dài của vectơ $\overrightarrow{BC}$.
		\end{enumerate}
	\end{btxanh}
	
	\begin{btvang}[]
		Cho hình lăng trụ tam giác $ABC.A'B'C'$. Sử dụng các đỉnh của hình lăng trụ làm điểm đầu và điểm cuối của các vectơ.
		\begin{enumerate}[label=\alph*)]
			\item Hãy kể tên các vectơ bằng nhau lần lượt bằng các vectơ $\overrightarrow{AA'}, \overrightarrow{BC}, \overrightarrow{A'C'}$.
			\item Hãy kể tên các vectơ luôn có độ dài bằng nhau và bằng độ dài của vectơ $\overrightarrow{AB}$.
		\end{enumerate}
	\end{btvang}
	
	\begin{btdo}[]
		Cho hình hộp chữ nhật $MNPQ.M'N'P'Q'$. Sử dụng các đỉnh của hình hộp chữ nhật làm điểm đầu và điểm cuối của các vectơ.
		\begin{enumerate}[label=\alph*)]
			\item Hãy kể tên các vectơ bằng nhau lần lượt bằng các vectơ $\overrightarrow{MN}, \overrightarrow{MP}, \overrightarrow{MQ}, \overrightarrow{MM'}$.
			\item Hãy kể tên các vectơ luôn có độ dài bằng nhau và bằng độ dài của vectơ $\overrightarrow{NP}$.
		\end{enumerate}
	\end{btdo}
	
	\subsection{Dạng 2: Chứng minh đẳng thức vector}	
	
	
	\begin{btxanh}[]
		\textbf{Câu 1:} Cho hình chóp $S.ABCD$ có đáy $ABCD$ là hình bình hành.
		\begin{enumerate}[label=\alph*)]
			\item Chứng minh rằng: $\overrightarrow{SA} + \overrightarrow{SC} = \overrightarrow{SB} + \overrightarrow{SD}$.
			\item Nếu đáy $ABCD$ là hình chữ nhật, chứng minh rằng: $\overrightarrow{SA}^2 + \overrightarrow{SC}^2 = \overrightarrow{SB}^2 + \overrightarrow{SD}^2$.
		\end{enumerate}
		
		\vspace{0.2cm}
		\textbf{Câu 2:} Cho hình hộp $ABCD.A'B'C'D'$. Chứng minh rằng:
		$$ \overrightarrow{AB} + \overrightarrow{B'C'} + \overrightarrow{CD} + \overrightarrow{D'A} = \vec{0} $$
		
		\vspace{0.2cm}
		\textbf{Câu 3:} Cho tứ diện $ABCD$. Gọi $M$ và $N$ lần lượt là trung điểm của các cạnh đối $AD$ và $BC$. Chứng minh rằng:
		$$ \overrightarrow{MN} = \frac{1}{2}(\overrightarrow{AB} + \overrightarrow{DC}) $$
	\end{btxanh}
	
	\begin{btvang}[]
		\textbf{Câu 1:} Cho hình chóp $S.XYZT$ có đáy $XYZT$ là hình bình hành.
		\begin{enumerate}[label=\alph*)]
			\item Chứng minh rằng: $\overrightarrow{SX} + \overrightarrow{SZ} = \overrightarrow{SY} + \overrightarrow{ST}$.
			\item Nếu đáy $XYZT$ là hình chữ nhật, chứng minh rằng: $\overrightarrow{SX}^2 + \overrightarrow{SZ}^2 = \overrightarrow{SY}^2 + \overrightarrow{ST}^2$.
		\end{enumerate}
		
		\vspace{0.2cm}
		\textbf{Câu 2:} Cho hình hộp $ABCD.A'B'C'D'$. Chứng minh rằng:
		$$ \overrightarrow{AB} + \overrightarrow{AD'} + \overrightarrow{C'C} + \overrightarrow{B'A'} = \overrightarrow{AD} $$
		
		\vspace{0.2cm}
		\textbf{Câu 3:} Cho tứ diện $S.ABC$. Gọi $G$ là trọng tâm của tam giác đáy $ABC$. Chứng minh rằng:
		$$ \overrightarrow{SA} + \overrightarrow{SB} + \overrightarrow{SC} = 3\overrightarrow{SG} $$
	\end{btvang}
	
	\begin{btdo}[]
		\textbf{Câu 1:} Cho hình chóp $S.ABCD$ có đáy là hình bình hành tâm $O$. Gọi $G$ là điểm thỏa mãn hệ thức: $\overrightarrow{GS} + \overrightarrow{GA} + \overrightarrow{GB} + \overrightarrow{GC} + \overrightarrow{GD} = \vec{0}$. Chứng minh rằng:
		$$ \overrightarrow{GS} = 4\overrightarrow{OG} $$
		
		\vspace{0.2cm}
		\textbf{Câu 2:} Cho hình hộp $ABCD.A'B'C'D'$. Chứng minh rằng:
		$$ \overrightarrow{AB} + \overrightarrow{AD} + \overrightarrow{AA'} + \overrightarrow{C'D'} + \overrightarrow{C'B} = \vec{0} $$
		
		\vspace{0.2cm}
		\textbf{Câu 3:} Cho hình hộp $ABCD.A'B'C'D'$ có tâm $O$. Gọi $M$ là tâm của hình bình hành đáy $ABCD$. Chứng minh rằng:
		$$ \overrightarrow{OM} = -\frac{1}{8}(\overrightarrow{AC'} + \overrightarrow{CA'} + \overrightarrow{BD'} + \overrightarrow{DB'}) $$
	\end{btdo}
	
	\subsection{Dạng 3: Phân tích vector}
	
	\begin{btxanh}[]
		\textbf{Câu 1:} Cho hình lăng trụ tam giác $ABC.A'B'C'$ có $\overrightarrow{AA'} = \vec{a}$, $\overrightarrow{AB} = \vec{b}$, $\overrightarrow{AC} = \vec{c}$. Hãy biểu diễn (phân tích) các vectơ $\overrightarrow{B'C}$ và $\overrightarrow{BC'}$ theo các vectơ $\vec{a}$, $\vec{b}$, $\vec{c}$.
		
		\vspace{0.2cm}
		\textbf{Câu 2:} Cho tứ diện $ABCD$.
		\begin{enumerate}[label=\alph*)]
			\item Gọi $I, J$ lần lượt là trung điểm của $AB$ và $CD$. Hãy biểu diễn vectơ $\overrightarrow{IJ}$ theo ba vectơ $\overrightarrow{AB}$, $\overrightarrow{AC}$ và $\overrightarrow{AD}$.
			\item Gọi $G$ là trọng tâm của tam giác $BCD$. Hãy biểu diễn vectơ $\overrightarrow{AG}$ theo ba vectơ $\overrightarrow{AB}$, $\overrightarrow{AC}$ và $\overrightarrow{AD}$.
		\end{enumerate}
		
		\vspace{0.2cm}
		\textbf{Câu 3:} Cho hình chóp $S.ABC$. Lấy hai điểm $M$ và $N$ sao cho $\overrightarrow{MS} = -2\overrightarrow{MA}$ và $\overrightarrow{NC} = -2\overrightarrow{NB}$. Chứng minh rằng:
		$$ \overrightarrow{MN} = \frac{1}{3}\overrightarrow{SC} + \frac{2}{3}\overrightarrow{AB} $$
	\end{btxanh}
	
	\begin{btvang}[]
		\textbf{Câu 1:} Cho tứ diện $ABCD$. Gọi $M$ và $N$ lần lượt là các điểm thuộc các cạnh $AD$ và $BC$ sao cho $AM = 2MD$ và $BC = 3NC$. Chứng minh rằng:
		$$ \overrightarrow{MN} = \frac{1}{3}\overrightarrow{AB} - \frac{2}{3}\overrightarrow{CD} $$
		
		\vspace{0.2cm}
		\textbf{Câu 2:} Cho tứ diện $ABCD$. Lấy các điểm $M$ và $N$ lần lượt thuộc các cạnh $AD$ và $BC$ sao cho $\overrightarrow{AM} = 3\overrightarrow{MD}$ và $\overrightarrow{NB} = -3\overrightarrow{NC}$. Đặt $\overrightarrow{AB} = \vec{a}$ và $\overrightarrow{CD} = \vec{b}$. Hãy biểu diễn vectơ $\overrightarrow{MN}$ theo hai vectơ $\vec{a}$ và $\vec{b}$.
		
		\vspace{0.2cm}
		\textbf{Câu 3:} Cho hình hộp $ABCD.A'B'C'D'$. Gọi $G$ và $G'$ lần lượt là trọng tâm của hai tam giác $BDA'$ và $CB'D'$. Chứng minh rằng ba điểm $A, G, C'$ thẳng hàng.
	\end{btvang}
	
	\begin{btdo}[]
		\textbf{Câu 1:} Cho hình hộp $ABCD.A'B'C'D'$. Đặt $\overrightarrow{BA} = \vec{a}$, $\overrightarrow{BB'} = \vec{b}$, $\overrightarrow{BC} = \vec{c}$. Gọi $M$ và $N$ lần lượt là hai điểm nằm trên các đường chéo $AC$ và $DC'$ sao cho $MB \parallel BD'$. Tính tỷ số độ dài $\frac{MN}{BD'}$.
		
		\vspace{0.2cm}
		\textbf{Câu 2:} Cho hình chóp $S.ABC$. Lấy các điểm $A', B', C'$ lần lượt thuộc các tia $SA, SB, SC$ sao cho $SA = a \cdot SA'$, $SB = b \cdot SB'$, $SC = c \cdot SC'$ (với $a, b, c$ là các số thực dương thay đổi). Chứng minh rằng mặt phẳng $(A'B'C')$ đi qua trọng tâm của tam giác $ABC$ khi và chỉ khi:
		$$ a + b + c = 3 $$
		
		\vspace{0.2cm}
		\textbf{Câu 3:} Cho tứ diện $ABCD$. Gọi $M$ và $N$ là các điểm lần lượt thuộc $AB$ và $CD$ sao cho $\overrightarrow{MA} = -2\overrightarrow{MB}$ và $\overrightarrow{ND} = -2\overrightarrow{NC}$. Các điểm $I, J, K$ lần lượt thuộc các đường thẳng $AD, MN, BC$ sao cho $\overrightarrow{IA} = k\overrightarrow{ID}$, $\overrightarrow{JM} = k\overrightarrow{JN}$, $\overrightarrow{KB} = k\overrightarrow{KC}$ (với $k \neq 1$). Chứng minh rằng ba điểm $I, J, K$ thẳng hàng.
	\end{btdo}
	
	\subsection{Dạng 4: Góc giữa hai vectơ và tích vô hướng}
	
	\begin{btxanh}[]
		\textbf{Câu 1:} Cho hình lập phương $ABCD.A'B'C'D'$. Hãy xác định góc giữa các cặp vectơ sau:
		\begin{enumerate}[label=\alph*)]
			\item $\overrightarrow{AB}$ và $\overrightarrow{D'D}$.
			\item $\overrightarrow{AB}$ và $\overrightarrow{A'C'}$.
		\end{enumerate}
		
		\vspace{0.2cm}
		\textbf{Câu 2:} Cho tứ diện đều $ABCD$ có $H$ là trung điểm của cạnh $AB$. Tính góc giữa các cặp vectơ sau:
		\begin{enumerate}[label=\alph*)]
			\item $\overrightarrow{AB}$ và $\overrightarrow{BC}$.
			\item $\overrightarrow{CH}$ và $\overrightarrow{AC}$.
		\end{enumerate}
		
		\vspace{0.2cm}
		\textbf{Câu 3:} Cho tứ diện $ABCD$ có $AB = AC = AD$ và $\widehat{BAC} = \widehat{BAD} = 60^\circ$. Chứng minh rằng:
		$$ \overrightarrow{AB} \cdot \overrightarrow{CD} = 0 $$
		Từ đó suy ra góc giữa cặp vectơ $\overrightarrow{AB}$ và $\overrightarrow{CD}$.
	\end{btxanh}
	
	\begin{btvang}[]
		\textbf{Câu 1:} Cho hình chóp $S.ABC$ có $SA = SB = SC$ và $\widehat{ASB} = \widehat{BSC} = \widehat{CSA}$. Hãy xác định góc giữa cặp vectơ $\overrightarrow{AB}$ và $\overrightarrow{SC}$.
		
		\vspace{0.2cm}
		\textbf{Câu 2:} Cho tứ diện $ABCD$ có $AB \perp AC$ và $AB \perp BD$. Gọi $P$ và $Q$ lần lượt là trung điểm của $AB$ và $CD$. Chứng minh rằng:
		$$ AB \perp PQ $$
		
		\vspace{0.2cm}
		\textbf{Câu 3:} Cho tứ diện $ABCD$ có $AB = AC = AD$ và $\widehat{BAC} = \widehat{BAD} = 60^\circ$, $\widehat{CAD} = 90^\circ$. Gọi $I$ và $J$ lần lượt là trung điểm của $AB$ và $CD$. Hãy xác định góc giữa cặp vectơ $\overrightarrow{AB}$ và $\overrightarrow{IJ}$.
	\end{btvang}
	
	\begin{btdo}[]
		\textbf{Câu 1:} Cho tứ diện $OABC$ có $OA, OB, OC$ đôi một vuông góc và $OA = OB = OC = a$. Gọi $M$ là trung điểm của $AB$. Tính góc giữa hai vectơ $\overrightarrow{OM}$ và $\overrightarrow{BC}$.
		
		\vspace{0.2cm}
		\textbf{Câu 2:} Cho hình hộp chữ nhật $ABCD.A'B'C'D'$ có $AB = a, AD = b, AA' = c$.
		\begin{enumerate}[label=\alph*)]
			\item Tính tích vô hướng của hai vectơ $\overrightarrow{AB}$ và $\overrightarrow{A'C'}$.
			\item Tính góc giữa hai vectơ này trong trường hợp đáy là hình vuông ($a = b$).
		\end{enumerate}
		
		\vspace{0.2cm}
		\textbf{Câu 3:} Cho tứ diện đều $ABCD$ cạnh $a$. Gọi $O$ là tâm đường tròn ngoại tiếp tam giác đáy $BCD$. Chứng minh rằng đường thẳng $AO$ vuông góc với $CD$, từ đó tính tích vô hướng:
		$$ \overrightarrow{AO} \cdot \overrightarrow{CD} $$
	\end{btdo}
	
	
	\section{Câu hỏi trắc nghiệm}
	
	\textbf{Câu 1:} Cho tứ diện $ABCD$. Hỏi có bao nhiêu vectơ khác vectơ $\vec{0}$ mà mỗi vectơ có điểm đầu và điểm cuối là hai đỉnh trong các đỉnh của tứ diện $ABCD$?
	\begin{multicols}{2}
		A. 12. \par
		B. 4. \par
		C. 10. \par
		D. 8.
	\end{multicols}
	
	\vspace{0.2cm}
	\noindent \textbf{Câu 2:} Cho hình chóp tứ giác $S.EFGH$. Hỏi có bao nhiêu vectơ khác vectơ $\vec{0}$ mà điểm đầu và điểm cuối thuộc tập hợp các đỉnh của hình chóp?
	\begin{multicols}{2}
		A. 20. \par
		B. 12. \par
		C. 16. \par
		D. 30.
	\end{multicols}
	
	\vspace{0.2cm}
	\noindent \textbf{Câu 3:} Cho lăng trụ tam giác $ABC.A'B'C'$. Số lượng vectơ khác vectơ $\vec{0}$ có điểm đầu và điểm cuối là các đỉnh của lăng trụ này là:
	\begin{multicols}{2}
		A. 30. \par
		B. 42. \par
		C. 24. \par
		D. 36.
	\end{multicols}
	
	\vspace{0.2cm}
	\noindent \textbf{Câu 4:} Cho hình bát diện đều $ABCDEF$. Số các vectơ khác vectơ $\vec{0}$ mà có điểm đầu và điểm cuối là hai đỉnh phân biệt của hình bát diện đều là:
	\begin{multicols}{2}
		A. 42. \par
		B. 30. \par
		C. 56. \par
		D. 20.
	\end{multicols}
	
	\noindent\textbf{Câu 5:} Cho hình hộp $ABCD.A'B'C'D'$. Gọi $I, J$ lần lượt là trung điểm của $AB'$ và $CD'$. Khẳng định nào dưới đây là đúng?
	\begin{multicols}{2}
		A. $\overrightarrow{AI} = \overrightarrow{CJ}$. \par
		B. $\overrightarrow{D'A'} = \overrightarrow{IJ}$. \par
		C. $\overrightarrow{BI} = \overrightarrow{D'J}$. \par
		D. $\overrightarrow{A'I} = \overrightarrow{JC}$.
	\end{multicols}
	
	\vspace{0.2cm}
	\noindent\textbf{Câu 6:} Cho hình lập phương $ABCD.A'B'C'D'$. Mệnh đề nào sau đây \textbf{sai}?
	\begin{multicols}{2}
		A. $\overrightarrow{AB} + \overrightarrow{AD} + \overrightarrow{AA'} = \overrightarrow{AC'}$. \par
		B. $\overrightarrow{AC} = \overrightarrow{AB} + \overrightarrow{AD}$. \par
		C. $|\overrightarrow{AB}| = |\overrightarrow{CD}|$. \par
		D. $\overrightarrow{AB} = \overrightarrow{CD}$.
	\end{multicols}
	
	\vspace{0.2cm}
	\noindent\textbf{Câu 7:} Cho tứ diện $ABCD$ có trọng tâm $G$. Mệnh đề nào sau đây \textbf{sai}?
	\begin{multicols}{2}
		A. $\overrightarrow{GA} + \overrightarrow{GB} + \overrightarrow{GC} + \overrightarrow{GD} = \vec{0}$. \par
		B. $\overrightarrow{OG} = \frac{1}{4}(\overrightarrow{OA} + \overrightarrow{OB} + \overrightarrow{OC} + \overrightarrow{OD})$. \par
		C. $\overrightarrow{AG} = \frac{2}{3}(\overrightarrow{AB} + \overrightarrow{AC} + \overrightarrow{AD})$. \par
		D. $\overrightarrow{AG} = \frac{1}{4}(\overrightarrow{AB} + \overrightarrow{AC} + \overrightarrow{AD})$.
	\end{multicols}
	
	\vspace{0.2cm}
	\noindent\textbf{Câu 8:} Cho tứ diện $ABCD$, gọi $I, J$ lần lượt là trung điểm của $AB$ và $CD$. Đẳng thức nào sau đây \textbf{sai}?
	\begin{multicols}{2}
		A. $\overrightarrow{IJ} = \frac{1}{2}(\overrightarrow{AC} + \overrightarrow{BD})$. \par
		B. $\overrightarrow{IJ} = \frac{1}{2}(\overrightarrow{AD} + \overrightarrow{BC})$. \par
		C. $\overrightarrow{IJ} = \frac{1}{2}(\overrightarrow{DC} + \overrightarrow{AD} + \overrightarrow{BD})$. \par
		D. $\overrightarrow{IJ} = \frac{1}{2}(\overrightarrow{AB} + \overrightarrow{CD})$.
	\end{multicols}
	
	\vspace{0.2cm}
	\noindent\textbf{Câu 9:} Cho tứ diện $ABCD$. Mệnh đề nào dưới đây là mệnh đề đúng?
	\begin{multicols}{2}
		A. $\overrightarrow{BC} + \overrightarrow{AB} = \overrightarrow{DA} - \overrightarrow{DC}$. \par
		B. $\overrightarrow{AC} - \overrightarrow{AD} = \overrightarrow{BD} - \overrightarrow{BC}$. \par
		C. $\overrightarrow{AB} - \overrightarrow{AC} = \overrightarrow{DB} - \overrightarrow{DC}$. \par
		D. $\overrightarrow{AB} - \overrightarrow{AD} = \overrightarrow{CD} + \overrightarrow{BC}$.
	\end{multicols}
	
	\vspace{0.2cm}
	\noindent\textbf{Câu 10:} Cho hình hộp $ABCD.A'B'C'D'$. Biểu thức nào sau đây bằng với vectơ $\overrightarrow{A'C}$?
	\begin{multicols}{2}
		A. $\overrightarrow{A'B'} + \overrightarrow{A'D'} + \overrightarrow{A'A}$. \par
		B. $\overrightarrow{A'B'} + \overrightarrow{A'D'} - \overrightarrow{A'A}$. \par
		C. $\overrightarrow{A'B'} - \overrightarrow{A'D'} + \overrightarrow{A'A}$. \par
		D. $\overrightarrow{A'B'} + \overrightarrow{A'D'} - 2\overrightarrow{A'A}$.
	\end{multicols}
	
	\vspace{0.2cm}
	\noindent\textbf{Câu 11:} Cho tứ diện $ABCD$. Gọi $M$ là trung điểm của cạnh $BC$. Mệnh đề nào sau đây đúng?
	\begin{multicols}{2}
		A. $\overrightarrow{AM} = \frac{1}{2}(\overrightarrow{AB} + \overrightarrow{AC})$. \par
		B. $\overrightarrow{AM} = \frac{1}{2}(\overrightarrow{AB} + \overrightarrow{AD})$. \par
		C. $\overrightarrow{AM} = \overrightarrow{AB} + \overrightarrow{AC}$. \par
		D. $\overrightarrow{AM} = \frac{1}{2}(\overrightarrow{AB} - \overrightarrow{AC})$.
	\end{multicols}
	
	\vspace{0.2cm}
	\noindent\textbf{Câu 12:} Cho hình chóp $S.ABC$. Gọi $G$ là trọng tâm tam giác $ABC$. Đẳng thức nào dưới đây là đúng?
	\begin{multicols}{2}
		A. $\overrightarrow{SA} + \overrightarrow{SB} + \overrightarrow{SC} = \overrightarrow{SG}$. \par
		B. $\overrightarrow{SA} + \overrightarrow{SB} + \overrightarrow{SC} = 3\overrightarrow{SG}$. \par
		C. $\overrightarrow{SA} + \overrightarrow{SB} + \overrightarrow{SC} = 2\overrightarrow{SG}$. \par
		D. $\overrightarrow{SA} + \overrightarrow{SB} + \overrightarrow{SC} = \vec{0}$.
	\end{multicols}
	
	\vspace{0.2cm}
	\noindent\textbf{Câu 13:} Cho hình lập phương $ABCD.A'B'C'D'$. Khẳng định nào sau đây là đúng về độ dài của vectơ?
	\begin{multicols}{2}
		A. $|\overrightarrow{AB} + \overrightarrow{AD} + \overrightarrow{AA'}| = A'C$. \par
		B. $|\overrightarrow{AB} + \overrightarrow{AD} + \overrightarrow{AA'}| = AC'$. \par
		C. $|\overrightarrow{AB} + \overrightarrow{AD} + \overrightarrow{AA'}| = 3AB$. \par
		D. $|\overrightarrow{AB} + \overrightarrow{AD} + \overrightarrow{AA'}| = A'B$.
	\end{multicols}
	
	\vspace{0.2cm}
	\noindent\textbf{Câu 14:} Cho tứ diện $ABCD$ có trọng tâm là $G$. Với điểm $M$ bất kì trong không gian, khẳng định nào sau đây luôn đúng?
	\begin{multicols}{2}
		A. $\overrightarrow{MA} + \overrightarrow{MB} + \overrightarrow{MC} + \overrightarrow{MD} = 4\overrightarrow{MG}$. \par
		B. $\overrightarrow{MA} + \overrightarrow{MB} + \overrightarrow{MC} + \overrightarrow{MD} = \overrightarrow{MG}$. \par
		C. $\overrightarrow{MA} + \overrightarrow{MB} + \overrightarrow{MC} + \overrightarrow{MD} = 2\overrightarrow{MG}$. \par
		D. $\overrightarrow{MA} + \overrightarrow{MB} + \overrightarrow{MC} + \overrightarrow{MD} = \vec{0}$.
	\end{multicols}
	
	\noindent\textbf{Câu 15:} Cho tứ diện $ABCD$. Gọi $M, N$ lần lượt là trung điểm của $AD$ và $BC$. Khẳng định nào sau đây \textbf{sai}?
	\begin{multicols}{2}
		A. $\overrightarrow{AB} + \overrightarrow{CD} = \overrightarrow{CB} + \overrightarrow{AD}$. \par
		B. $2\overrightarrow{MN} = \overrightarrow{AB} + \overrightarrow{DC}$. \par
		C. $\overrightarrow{AD} + 2\overrightarrow{MN} = \overrightarrow{AB} + \overrightarrow{AC}$. \par
		D. $2\overrightarrow{MN} = \overrightarrow{AB} + \overrightarrow{AC} + \overrightarrow{AD}$.
	\end{multicols}
	
	\vspace{0.2cm}
	\noindent\textbf{Câu 16:} Cho hình chóp $S.ABCD$ có đáy $ABCD$ là hình bình hành. Khẳng định nào sau đây \textbf{đúng}?
	\begin{multicols}{2}
		A. $\overrightarrow{SA} + \overrightarrow{SD} = \overrightarrow{SB} + \overrightarrow{SC}$. \par
		B. $\overrightarrow{SA} + \overrightarrow{SB} + \overrightarrow{SC} + \overrightarrow{SD} = \vec{0}$. \par
		C. $\overrightarrow{SA} + \overrightarrow{SC} = \overrightarrow{SB} + \overrightarrow{SD}$. \par
		D. $\overrightarrow{SA} + \overrightarrow{SB} = \overrightarrow{SC} + \overrightarrow{SD}$.
	\end{multicols}
	
	\vspace{0.2cm}
	\noindent\textbf{Câu 17:} Cho tứ diện $ABCD$. Gọi $I, J$ lần lượt là trung điểm của $AB$ và $CD$, $G$ là trung điểm của $IJ$. Đẳng thức nào sau đây đúng?
	\begin{multicols}{2}
		A. $\overrightarrow{GA} + \overrightarrow{GB} + \overrightarrow{GC} + \overrightarrow{GD} = \vec{0}$. \par
		B. $\overrightarrow{GA} + \overrightarrow{GB} + \overrightarrow{GC} + \overrightarrow{GD} = 2\overrightarrow{IJ}$. \par
		C. $\overrightarrow{GA} + \overrightarrow{GB} + \overrightarrow{GC} + \overrightarrow{GD} = \overrightarrow{JI}$. \par
		D. $\overrightarrow{GA} + \overrightarrow{GB} + \overrightarrow{GC} + \overrightarrow{GD} = -2\overrightarrow{JI}$.
	\end{multicols}
	
	\vspace{0.2cm}
	\noindent\textbf{Câu 18:} Cho hình lăng trụ tam giác $ABC.A'B'C'$. Đặt $\overrightarrow{AA'} = \vec{a}$, $\overrightarrow{AB} = \vec{b}$, $\overrightarrow{AC} = \vec{c}$, $\overrightarrow{BC} = \vec{d}$. Biểu thức vectơ nào sau đây đúng?
	\begin{multicols}{2}
		A. $\vec{a} + \vec{b} + \vec{c} = \vec{d}$. \par
		B. $\vec{a} = \vec{b} + \vec{c}$. \par
		C. $\vec{a} + \vec{b} + \vec{c} + \vec{d} = \vec{0}$. \par
		D. $\vec{b} - \vec{c} + \vec{d} = \vec{0}$.
	\end{multicols}
	
	\vspace{0.2cm}
	\noindent\textbf{Câu 19:} Cho hình hộp chữ nhật $ABCD.A'B'C'D'$. Đặt $\overrightarrow{AB} = \vec{x}$, $\overrightarrow{AD} = \vec{y}$, $\overrightarrow{AA'} = \vec{z}$. Vectơ biểu diễn đường chéo $\overrightarrow{BD'}$ theo $\vec{x}$, $\vec{y}$, $\vec{z}$ là:
	\begin{multicols}{2}
		A. $\overrightarrow{BD'} = -\vec{x} + \vec{y} + \vec{z}$. \par
		B. $\overrightarrow{BD'} = \vec{x} - \vec{y} + \vec{z}$. \par
		C. $\overrightarrow{BD'} = \vec{x} + \vec{y} - \vec{z}$. \par
		D. $\overrightarrow{BD'} = \vec{x} + \vec{y} + \vec{z}$.
	\end{multicols}
	
	\vspace{0.2cm}
	\noindent\textbf{Câu 20:} Cho hình lăng trụ tam giác $ABC.A'B'C'$. Gọi $M$ là trung điểm của $B'C'$. Đẳng thức nào sau đây là đúng?
	\begin{multicols}{2}
		A. $\overrightarrow{AM} = \overrightarrow{AA'} + \frac{1}{2}(\overrightarrow{AB} + \overrightarrow{AC})$. \par
		B. $\overrightarrow{AM} = \overrightarrow{AA'} + \overrightarrow{AB} + \overrightarrow{AC}$. \par
		C. $\overrightarrow{AM} = \frac{1}{2}\overrightarrow{AA'} + \overrightarrow{AB} + \overrightarrow{AC}$. \par
		D. $\overrightarrow{AM} = \overrightarrow{AA'} + \frac{1}{2}(\overrightarrow{AB} - \overrightarrow{AC})$.
	\end{multicols}
	
	\vspace{0.2cm}
	\noindent\textbf{Câu 21:} Cho tứ diện đều $ABCD$ có cạnh $a$. Gọi $H$ là trực tâm tam giác đáy $BCD$. Đẳng thức nào sau đây là đúng?
	\begin{multicols}{2}
		A. $\overrightarrow{AH} = \frac{1}{3}(\overrightarrow{AB} + \overrightarrow{AC} + \overrightarrow{AD})$. \par
		B. $\overrightarrow{AH} = \frac{1}{2}(\overrightarrow{AB} + \overrightarrow{AC} + \overrightarrow{AD})$. \par
		C. $\overrightarrow{AH} = \overrightarrow{AB} + \overrightarrow{AC} + \overrightarrow{AD}$. \par
		D. $\overrightarrow{AH} = \frac{1}{4}(\overrightarrow{AB} + \overrightarrow{AC} + \overrightarrow{AD})$.
	\end{multicols}
	
	\vspace{0.2cm}
	\noindent\textbf{Câu 22:} Cho hình chóp $S.ABCD$ có đáy $ABCD$ là hình thang có hai đáy $AB$ và $CD$ sao cho $AB = 2CD$. Gọi $M$ là trung điểm của $AB$. Khẳng định nào sau đây là đúng?
	\begin{multicols}{2}
		A. $\overrightarrow{DM} = \overrightarrow{CB}$. \par
		B. $\overrightarrow{DM} = \overrightarrow{BC}$. \par
		C. $\overrightarrow{DM} = \overrightarrow{AD}$. \par
		D. $\overrightarrow{DM} = \overrightarrow{CD}$.
	\end{multicols}
	
	\vspace{0.2cm}
	\noindent\textbf{Câu 23:} Cho hình hộp $ABCD.A'B'C'D'$ tâm $O$. Đẳng thức nào sau đây \textbf{sai}?
	\begin{multicols}{2}
		A. $\overrightarrow{OA} + \overrightarrow{OC'} = \vec{0}$. \par
		B. $\overrightarrow{OB} + \overrightarrow{OD'} = \vec{0}$. \par
		C. $\overrightarrow{OA} + \overrightarrow{OB} + \overrightarrow{OC} + \overrightarrow{OD} = \vec{0}$. \par
		D. $\overrightarrow{OA} + \overrightarrow{OB'} + \overrightarrow{OC} + \overrightarrow{OD'} = \vec{0}$.
	\end{multicols}
	
	\vspace{0.2cm}
	\noindent\textbf{Câu 24:} Cho tứ diện $ABCD$. Gọi $G_1, G_2, G_3, G_4$ lần lượt là trọng tâm của các tam giác $BCD, CDA, DAB, ABC$. Gọi $G$ là trọng tâm của tứ diện $ABCD$. Khẳng định nào sau đây đúng?
	\begin{multicols}{2}
		A. $\overrightarrow{GG_1} + \overrightarrow{GG_2} + \overrightarrow{GG_3} + \overrightarrow{GG_4} = \vec{0}$. \par
		B. $\overrightarrow{AG_1} + \overrightarrow{BG_2} + \overrightarrow{CG_3} + \overrightarrow{DG_4} = \vec{0}$. \par
		C. $\overrightarrow{AG_1} = \overrightarrow{BG_2}$. \par
		D. Cả A và B đều đúng.
	\end{multicols}
	
	\noindent \textbf{Câu 25:} Cho hình chóp $S.XYZT$ có đáy $XYZT$ là hình bình hành. Đặt $\overrightarrow{SX} = \vec{u}$, $\overrightarrow{SY} = \vec{v}$, $\overrightarrow{SZ} = \vec{w}$, $\overrightarrow{ST} = \vec{t}$. Khẳng định nào sau đây đúng?
	\begin{multicols}{2}
		A. $\vec{u} + \vec{w} = \vec{v} + \vec{t}$. \par
		B. $\vec{u} + \vec{v} = \vec{w} + \vec{t}$. \par
		C. $\vec{u} + \vec{t} = \vec{v} + \vec{w}$. \par
		D. $\vec{u} + \vec{v} + \vec{w} + \vec{t} = \vec{0}$.
	\end{multicols}
	
	\vspace{0.2cm}
	\noindent \textbf{Câu 26:} Cho hình chóp $S.ABC$. Gọi $G$ là trọng tâm của tam giác $ABC$. Đẳng thức nào sau đây đúng?
	\begin{multicols}{2}
		A. $\overrightarrow{SG} = \overrightarrow{SA} + \overrightarrow{SB} + \overrightarrow{SC}$. \par
		B. $\overrightarrow{SG} = \frac{1}{3}(\overrightarrow{SA} + \overrightarrow{SB} + \overrightarrow{SC})$. \par
		C. $\overrightarrow{SG} = 3(\overrightarrow{SA} + \overrightarrow{SB} + \overrightarrow{SC})$. \par
		D. $\overrightarrow{SG} = \frac{1}{2}(\overrightarrow{SA} + \overrightarrow{SB} + \overrightarrow{SC})$.
	\end{multicols}
	
	\vspace{0.2cm}
	\noindent \textbf{Câu 27:} Cho hình chóp $S.ABCD$ có đáy $ABCD$ là hình bình hành. Trong các đẳng thức sau, đẳng thức nào luôn đúng?
	\begin{multicols}{2}
		A. $\overrightarrow{AB} + \overrightarrow{BC} + \overrightarrow{CD} + \overrightarrow{DA} = \vec{0}$. \par
		B. $\overrightarrow{SA} + \overrightarrow{SB} = \overrightarrow{SC} + \overrightarrow{SD}$. \par
		C. $\overrightarrow{AB} + \overrightarrow{AD} = \overrightarrow{BD}$. \par
		D. $\overrightarrow{SA} + \overrightarrow{SC} + \overrightarrow{SB} + \overrightarrow{SD} = \vec{0}$.
	\end{multicols}
	
	\vspace{0.2cm}
	\noindent \textbf{Câu 28:} Cho hình lập phương $ABCD.A'B'C'D'$. Thực hiện phép toán: $\vec{v} = \overrightarrow{A'B'} + \overrightarrow{A'D'} + \overrightarrow{A'A}$. Khẳng định nào dưới đây đúng?
	\begin{multicols}{2}
		A. $\vec{v} = \overrightarrow{A'C}$. \par
		B. $\vec{v} = \overrightarrow{A'C'}$. \par
		C. $\vec{v} = \overrightarrow{CA'}$. \par
		D. $\vec{v} = \overrightarrow{C'A'}$.
	\end{multicols}
	
	\vspace{0.2cm}
	\noindent \textbf{Câu 29:} Cho hình chóp $S.ABCD$ có đáy $ABCD$ là hình bình hành tâm $O$. Đẳng thức nào sau đây luôn đúng?
	\begin{multicols}{2}
		A. $\overrightarrow{SA} + \overrightarrow{SB} + \overrightarrow{SC} + \overrightarrow{SD} = 4\overrightarrow{SO}$. \par
		B. $\overrightarrow{SA} + \overrightarrow{SB} + \overrightarrow{SC} + \overrightarrow{SD} = 2\overrightarrow{SO}$. \par
		C. $\overrightarrow{SA} + \overrightarrow{SB} + \overrightarrow{SC} + \overrightarrow{SD} = \overrightarrow{SO}$. \par
		D. $\overrightarrow{SA} + \overrightarrow{SB} + \overrightarrow{SC} + \overrightarrow{SD} = \vec{0}$.
	\end{multicols}
	
	\noindent \textbf{Câu 30:} Cho hình lăng trụ tam giác $ABC.A'B'C'$. Đặt $\overrightarrow{AB} = \vec{a}$, $\overrightarrow{AA'} = \vec{b}$, $\overrightarrow{AC} = \vec{c}$. Khẳng định nào sau đây đúng?
	\begin{multicols}{2}
		A. $\overrightarrow{B'C} = \vec{a} + \vec{b} - \vec{c}$. \par
		B. $\overrightarrow{B'C} = -\vec{a} + \vec{b} - \vec{c}$. \par
		C. $\overrightarrow{B'C} = -\vec{a} - \vec{b} + \vec{c}$. \par
		D. $\overrightarrow{B'C} = -\vec{a} + \vec{b} + \vec{c}$.
	\end{multicols}
	
	\vspace{0.2cm}
	\noindent \textbf{Câu 31:} Cho tam giác $ABC$ có $AB = 3, AC = 4$, gọi $AD$ là đường phân giác trong của góc $A$. Mệnh đề nào sau đây đúng?
	\begin{multicols}{2}
		A. $\overrightarrow{AD} = \frac{4}{7}\overrightarrow{AB} + \frac{3}{7}\overrightarrow{AC}$. \par
		B. $\overrightarrow{AD} = \frac{4}{7}\overrightarrow{AB} - \frac{3}{7}\overrightarrow{AC}$. \par
		C. $\overrightarrow{AD} = -\frac{4}{7}\overrightarrow{AB} + \frac{3}{7}\overrightarrow{AC}$. \par
		D. $\overrightarrow{AD} = \frac{3}{7}\overrightarrow{AB} + \frac{4}{7}\overrightarrow{AC}$.
	\end{multicols}
	
	\vspace{0.2cm}
	\noindent \textbf{Câu 32:} Cho hình lăng trụ tam giác $ABC.A'B'C'$, gọi $M$ là trung điểm của cạnh bên $BB'$. Đặt $\overrightarrow{CA} = \vec{a}$, $\overrightarrow{CB} = \vec{b}$, $\overrightarrow{CC'} = \vec{c}$. Khẳng định nào sau đây là đúng?
	\begin{multicols}{2}
		A. $\overrightarrow{AM} = -\vec{a} + \vec{b} + \frac{1}{2}\vec{c}$. \par
		B. $\overrightarrow{AM} = \vec{a} - \frac{1}{2}\vec{b} + \vec{c}$. \par
		C. $\overrightarrow{AM} = -\frac{1}{2}\vec{a} + \vec{b} + \vec{c}$. \par
		D. $\overrightarrow{AM} = \vec{a} + \frac{1}{2}\vec{b} - \vec{c}$.
	\end{multicols}
	
	\vspace{0.2cm}
	\noindent \textbf{Câu 33:} Cho tứ diện $ABCD$. Gọi $M$ và $P$ lần lượt là trung điểm của $BC$ và $AD$. Đặt $\overrightarrow{AB} = \vec{b}$, $\overrightarrow{AC} = \vec{c}$, $\overrightarrow{AD} = \vec{d}$. Khẳng định nào sau đây đúng?
	\begin{multicols}{2}
		A. $\overrightarrow{MP} = \frac{1}{2}(\vec{d} + \vec{b} - \vec{c})$. \par
		B. $\overrightarrow{MP} = \frac{1}{2}(\vec{d} - \vec{b} - \vec{c})$. \par
		C. $\overrightarrow{MP} = \frac{1}{2}(\vec{c} + \vec{d} + \vec{b})$. \par
		D. $\overrightarrow{MP} = \frac{1}{2}(\vec{c} + \vec{b} - \vec{d})$.
	\end{multicols}
	
	\vspace{0.2cm}
	\noindent \textbf{Câu 34:} Cho tứ diện $ABCD$. Gọi $G$ là trọng tâm của tam giác $BCD$. Đặt $\overrightarrow{AB} = \vec{b}$, $\overrightarrow{AC} = \vec{c}$, $\overrightarrow{AD} = \vec{d}$. Khẳng định nào sau đây đúng?
	\begin{multicols}{2}
		A. $\overrightarrow{AG} = \frac{1}{3}(\vec{b} + \vec{c} + \vec{d})$. \par
		B. $\overrightarrow{AG} = \frac{1}{4}(\vec{b} + \vec{c} + \vec{d})$. \par
		C. $\overrightarrow{AG} = \vec{b} + \vec{c} + \vec{d}$. \par
		D. $\overrightarrow{AG} = \frac{1}{2}(\vec{b} + \vec{c} + \vec{d})$.
	\end{multicols}
	
	\vspace{0.2cm}
	\noindent \textbf{Câu 35:} Cho hình lăng trụ tam giác $ABC.A'B'C'$. Đặt $\overrightarrow{AA'} = \vec{a}$, $\overrightarrow{AB} = \vec{b}$, $\overrightarrow{AC} = \vec{c}$. Vectơ biểu diễn đường chéo $\overrightarrow{A'C}$ theo các vectơ $\vec{a}, \vec{b}, \vec{c}$ là:
	\begin{multicols}{2}
		A. $\overrightarrow{A'C} = -\vec{a} + \vec{c}$. \par
		B. $\overrightarrow{A'C} = -\vec{a} + \vec{b} + \vec{c}$. \par
		C. $\overrightarrow{A'C} = \vec{a} - \vec{c}$. \par
		D. $\overrightarrow{A'C} = \vec{a} + \vec{b} - \vec{c}$.
	\end{multicols}
	
	\noindent \textbf{Câu 36:} Cho tứ diện $ABCD$. Gọi $M$ và $P$ lần lượt là trung điểm của $AB$ và $CD$. Đặt $\overrightarrow{AB} = \vec{b}$, $\overrightarrow{AC} = \vec{c}$, $\overrightarrow{AD} = \vec{d}$. Khẳng định nào sau đây đúng?
	\begin{multicols}{2}
		A. $\overrightarrow{MP} = \frac{1}{2}(\vec{c} + \vec{b} - \vec{d})$. \par
		B. $\overrightarrow{MP} = \frac{1}{2}(\vec{c} + \vec{d} - \vec{b})$. \par
		C. $\overrightarrow{MP} = \frac{1}{2}(\vec{c} + \vec{d} + \vec{b})$. \par
		D. $\overrightarrow{MP} = \frac{1}{2}(\vec{d} + \vec{b} - \vec{c})$.
	\end{multicols}
	
	\vspace{0.2cm}
	\noindent \textbf{Câu 37:} Cho tứ diện $ABCD$. Đặt $\overrightarrow{AB} = \vec{a}$, $\overrightarrow{AC} = \vec{b}$, $\overrightarrow{AD} = \vec{c}$, gọi $M$ là trung điểm của $BC$. Trong các khẳng định sau, khẳng định nào đúng?
	\begin{multicols}{2}
		A. $\overrightarrow{DM} = \frac{1}{2}(\vec{a} + 2\vec{b} - \vec{c})$. \par
		B. $\overrightarrow{DM} = \frac{1}{2}(-2\vec{a} + \vec{b} + \vec{c})$. \par
		C. $\overrightarrow{DM} = \frac{1}{2}(\vec{a} - 2\vec{b} + \vec{c})$. \par
		D. $\overrightarrow{DM} = \frac{1}{2}(\vec{a} + \vec{b} - 2\vec{c})$.
	\end{multicols}
	
	\vspace{0.2cm}
	\noindent \textbf{Câu 38:} Cho tứ diện đều $ABCD$, $M$ là trung điểm của cạnh $AB$ và $G$ là trọng tâm của tam giác $BCD$. Đặt $\overrightarrow{AB} = \vec{b}$, $\overrightarrow{AC} = \vec{c}$, $\overrightarrow{AD} = \vec{d}$. Phân tích vectơ $\overrightarrow{MG}$ theo các vectơ $\vec{b}, \vec{c}, \vec{d}$ ta được:
	\begin{multicols}{2}
		A. $\overrightarrow{MG} = -\frac{1}{6}\vec{b} + \frac{1}{3}\vec{c} + \frac{1}{3}\vec{d}$. \par
		B. $\overrightarrow{MG} = \frac{1}{6}\vec{b} + \frac{1}{3}\vec{c} + \frac{1}{3}\vec{d}$. \par
		C. $\overrightarrow{MG} = -\frac{1}{6}\vec{b} - \frac{1}{3}\vec{c} + \frac{1}{3}\vec{d}$. \par
		D. $\overrightarrow{MG} = -\frac{1}{6}\vec{b} - \frac{1}{3}\vec{c} - \frac{1}{3}\vec{d}$.
	\end{multicols}
	
	\vspace{0.2cm}
	\noindent \textbf{Câu 39:} Cho tứ diện $ABCD$. Đặt $\overrightarrow{AB} = \vec{a}$, $\overrightarrow{AC} = \vec{b}$, $\overrightarrow{AD} = \vec{c}$, gọi $N$ là trung điểm của $AD$. Khẳng định nào sau đây là đúng?
	\begin{multicols}{2}
		A. $\overrightarrow{BN} = \vec{c} - \frac{1}{2}\vec{a}$. \par
		B. $\overrightarrow{CN} = \frac{1}{2}\vec{c} - \vec{b}$. \par
		C. $\overrightarrow{CN} = \vec{c} - \frac{1}{2}\vec{b}$. \par
		D. $\overrightarrow{BN} = \vec{a} - \frac{1}{2}\vec{c}$.
	\end{multicols}
	
	\vspace{0.2cm}
	\noindent \textbf{Câu 40:} Cho tứ diện $ABCD$. Gọi $I$ và $J$ lần lượt là trung điểm của $AC$ và $BD$. Đặt $\overrightarrow{AB} = \vec{b}$, $\overrightarrow{AC} = \vec{c}$, $\overrightarrow{AD} = \vec{d}$. Khẳng định nào sau đây đúng?
	\begin{multicols}{2}
		A. $\overrightarrow{IJ} = \frac{1}{2}(\vec{b} + \vec{d} - \vec{c})$. \par
		B. $\overrightarrow{IJ} = \frac{1}{2}(\vec{d} - \vec{b} - \vec{c})$. \par
		C. $\overrightarrow{IJ} = \frac{1}{2}(\vec{c} + \vec{d} - \vec{b})$. \par
		D. $\overrightarrow{IJ} = \frac{1}{2}(\vec{b} + \vec{c} - \vec{d})$.
	\end{multicols}
	
	\vspace{0.2cm}
	\noindent \textbf{Câu 41:} Cho tứ diện đều $ABCD$, $M$ là trung điểm của cạnh $AB$ và $G'$ là trọng tâm của tam giác $ACD$. Đặt $\overrightarrow{AB} = \vec{b}$, $\overrightarrow{AC} = \vec{c}$, $\overrightarrow{AD} = \vec{d}$. Phân tích vectơ $\overrightarrow{MG'}$ theo các vectơ $\vec{b}, \vec{c}, \vec{d}$ ta được:
	\begin{multicols}{2}
		A. $\overrightarrow{MG'} = -\frac{1}{2}\vec{b} + \frac{1}{3}\vec{c} + \frac{1}{3}\vec{d}$. \par
		B. $\overrightarrow{MG'} = \frac{1}{2}\vec{b} + \frac{1}{3}\vec{c} + \frac{1}{3}\vec{d}$. \par
		C. $\overrightarrow{MG'} = -\frac{1}{3}\vec{b} + \frac{1}{3}\vec{c} + \frac{1}{3}\vec{d}$. \par
		D. $\overrightarrow{MG'} = -\frac{1}{6}\vec{b} + \frac{1}{3}\vec{c} + \frac{1}{3}\vec{d}$.
	\end{multicols}
	% --- BÀI TẬP ĐÚNG/SAI ---
	\section{Câu hỏi đúng sai}
	
	\noindent \textbf{Câu 1:} Cho hình chóp đều $S.ABC$ có cạnh đáy bằng $a$ và cạnh bên bằng $2a$. Xét tính đúng sai của các mệnh đề sau:
	\begin{enumerate}[label=\alph*)]
		\item Có 4 vectơ có điểm đầu là $S$ và điểm cuối là một trong các đỉnh còn lại của tứ diện.
		\item Trong số các vectơ ở câu a, có 2 vectơ có giá nằm trong mặt phẳng $(ABC)$.
		\item Gọi $H$ là trung điểm của $AB$. Độ dài của vectơ $\overrightarrow{SH}$ bằng $\frac{a\sqrt{3}}{2}$.
		\item Gọi $G$ là trọng tâm của $\triangle ABC$. Độ dài của vectơ $\overrightarrow{SG}$ bằng $\frac{a\sqrt{3}}{3}$.
	\end{enumerate}
	
	\vspace{0.2cm}
	\noindent \textbf{Câu 2:} Cho hình hộp thoi $ABCD.A'B'C'D'$. Xét tính đúng sai của các mệnh đề sau:
	\begin{enumerate}[label=\alph*)]
		\item Trong các vectơ $\overrightarrow{CB}, \overrightarrow{C'B'}, \overrightarrow{B'C'}, \overrightarrow{AD}, \overrightarrow{DA}, \overrightarrow{A'D'}, \overrightarrow{D'A'}, \overrightarrow{AC}$ có 7 vectơ cùng phương với $\overrightarrow{BC}$.
		\item Hai vectơ $\overrightarrow{AC}$ và $\overrightarrow{A'C'}$ cùng hướng.
		\item Trong các vectơ $\overrightarrow{CB}, \overrightarrow{C'B'}, \overrightarrow{B'C'}, \overrightarrow{AD}, \overrightarrow{DA}, \overrightarrow{A'D'}, \overrightarrow{D'A'}, \overrightarrow{AC}$ có 6 vectơ bằng $\overrightarrow{BC}$.
		\item Nếu $A'B \perp (ADC'B')$ thì $|\overrightarrow{A'C}| = |\overrightarrow{BD'}|$.
	\end{enumerate}
	
	\vspace{0.2cm}
	\noindent \textbf{Câu 3:} Cho hình chóp đều $S.ABCD$ có tất cả các cạnh bằng 1 và $O$ là tâm của đáy. Xét tính đúng sai của các mệnh đề sau:
	\begin{enumerate}[label=\alph*)]
		\item Trong các vectơ $\overrightarrow{SA}, \overrightarrow{SB}, \overrightarrow{SC}, \overrightarrow{SD}$ có 4 vectơ cùng hướng với $\overrightarrow{SO}$.
		\item Số vectơ khác $\vec{0}$ cùng phương với $\overrightarrow{CD}$ và có điểm đầu, điểm cuối thuộc tập $\{S, A, B, O\}$ là 2.
		\item Độ dài của vectơ $\overrightarrow{SO}$ bằng $\frac{1}{4}$.
		\item Nếu $\overrightarrow{SM} = \overrightarrow{CA}$ thì $\widehat{DSM} = 120^\circ$.
	\end{enumerate}
	
	\noindent \textbf{Câu 4:} Cho hình lập phương $ABCD.A'B'C'D'$, có độ dài các cạnh bằng $1$. Xét tính đúng sai của các mệnh đề sau:
	\begin{enumerate}[label=\alph*)]
		\item $|\overrightarrow{BC} + \overrightarrow{DD'}| = \sqrt{2}$.
		\item $|\overrightarrow{AB} + \overrightarrow{AD} + \overrightarrow{AA'}| = 3$.
		\item $|\overrightarrow{AB} + \overrightarrow{A'A}| = \sqrt{3}$.
		\item $|\overrightarrow{A'C'} + \overrightarrow{C'C}| = \sqrt{3}$.
	\end{enumerate}
	
	\vspace{0.2cm}
	\noindent \textbf{Câu 5:} Cho hình lập phương $ABCD.A'B'C'D'$ có cạnh bằng $\sqrt{2}$. Xét tính đúng sai của các mệnh đề sau:
	\begin{enumerate}[label=\alph*)]
		\item $|\overrightarrow{AB} + \overrightarrow{CD}| = 2\sqrt{2}$.
		\item $|\overrightarrow{AB} + \overrightarrow{D'C'} - \overrightarrow{A'B'}| = \sqrt{2}$.
		\item $|\overrightarrow{AB} + \overrightarrow{A'D'}| = 4$.
		\item $|\overrightarrow{AB} + \overrightarrow{AD} + \overrightarrow{CC'}| = 3\sqrt{2}$.
	\end{enumerate}
	
	\vspace{0.2cm}
	\noindent \textbf{Câu 6:} Cho tứ diện $ABCD$. Xét tính đúng sai của các mệnh đề sau:
	\begin{enumerate}[label=\alph*)]
		\item $\overrightarrow{AC} + \overrightarrow{BD} = \overrightarrow{AD} + \overrightarrow{BC}$.
		\item $|\overrightarrow{AB} + \overrightarrow{BD}| = |\overrightarrow{AD}|$.
		\item $\overrightarrow{AB} + \overrightarrow{CD} = \overrightarrow{AD} + \overrightarrow{CB}$.
		\item $\overrightarrow{AB} + \overrightarrow{AC} + \overrightarrow{AD} = 3\overrightarrow{AG}$ với $G$ là trọng tâm tam giác $BCD$.
	\end{enumerate}
	
	\vspace{0.2cm}
	\noindent \textbf{Câu 7:} Cho lăng trụ $ABC.A'B'C'$, $O$ là giao điểm của $BC'$ và $B'C$. Xét tính đúng sai của các mệnh đề sau:
	\begin{enumerate}[label=\alph*)]
		\item $\overrightarrow{AB} + \overrightarrow{B'C'} - \overrightarrow{A'B} = \overrightarrow{BC'}$.
		\item $\overrightarrow{A'B} + \overrightarrow{AC} = \overrightarrow{A'C} + \overrightarrow{AB}$.
		\item $\overrightarrow{AC'} - \overrightarrow{BC'} = \overrightarrow{AC} - \overrightarrow{B'C'}$.
		\item Nếu $\overrightarrow{AA'} + \overrightarrow{AB} + \overrightarrow{AC} = \overrightarrow{AG}$ thì $G \equiv O$.
	\end{enumerate}
	
	\vspace{0.2cm}
	\noindent \textbf{Câu 8:} Cho hình hộp chữ nhật $ABCD.A'B'C'D'$ có $AB = 3, AD = 4, AA' = 5$. Xét tính đúng sai của các mệnh đề sau:
	\begin{enumerate}[label=\alph*)]
		\item $|\overrightarrow{AB} + \overrightarrow{AD}| = 5$.
		\item $|\overrightarrow{AB} + \overrightarrow{AA'}| = \sqrt{34}$.
		\item $|\overrightarrow{AB} + \overrightarrow{AD} + \overrightarrow{AA'}| = 5\sqrt{2}$.
		\item $|\overrightarrow{AB'} + \overrightarrow{C'D}| = 0$.
	\end{enumerate}
	
	\vspace{0.2cm}
	\noindent \textbf{Câu 9:} Cho tứ diện $ABCD$ có trọng tâm $G$. Xét tính đúng sai của các mệnh đề sau:
	\begin{enumerate}[label=\alph*)]
		\item $\overrightarrow{GA} + \overrightarrow{GB} + \overrightarrow{GC} + \overrightarrow{GD} = \vec{0}$.
		\item $\overrightarrow{OG} = \frac{1}{4}(\overrightarrow{OA} + \overrightarrow{OB} + \overrightarrow{OC} + \overrightarrow{OD})$ với mọi điểm $O$.
		\item $\overrightarrow{AG} = \frac{1}{4}(\overrightarrow{AB} + \overrightarrow{AC} + \overrightarrow{AD})$.
		\item $\overrightarrow{GB} + \overrightarrow{GC} = 2\overrightarrow{GM}$ với $M$ là trung điểm của $AD$.
	\end{enumerate}
	
	\vspace{0.2cm}
	\noindent \textbf{Câu 10:} Cho hình chóp $S.ABCD$ có đáy $ABCD$ là hình bình hành tâm $O$. Xét tính đúng sai của các mệnh đề sau:
	\begin{enumerate}[label=\alph*)]
		\item $\overrightarrow{SA} + \overrightarrow{SC} = \overrightarrow{SB} + \overrightarrow{SD}$.
		\item $\overrightarrow{SA} + \overrightarrow{SB} + \overrightarrow{SC} + \overrightarrow{SD} = 4\overrightarrow{SO}$.
		\item Nếu $\overrightarrow{SM} = \overrightarrow{AB}$ thì $AB \parallel SM$.
		\item Nếu $\overrightarrow{SA} + \overrightarrow{SB} + \overrightarrow{SC} + \overrightarrow{SD} = \vec{0}$ thì $S \equiv O$.
	\end{enumerate}
	
	\noindent \textbf{Câu 11:} Cho tứ diện $ABCD$. Các điểm $M, N, I$ lần lượt là trung điểm của $AB, CD, MN$ và $G$ là trọng tâm tam giác $BCD$. Xét tính đúng sai của các mệnh đề sau:
	\begin{enumerate}[label=\alph*)]
		\item $\overrightarrow{MC} + \overrightarrow{MD} = 4\overrightarrow{MN}$.
		\item $\overrightarrow{AD} + \overrightarrow{BC} = 2\overrightarrow{MN}$.
		\item $\overrightarrow{IB} + \overrightarrow{IC} + \overrightarrow{ID} = 3\overrightarrow{IG}$.
		\item $2\overrightarrow{IG} + \overrightarrow{IA} = \vec{0}$.
	\end{enumerate}
	
	\vspace{0.2cm}
	\noindent \textbf{Câu 12:} Cho hình lập phương $ABCD.A'B'C'D'$. Xét tính đúng sai của các mệnh đề sau:
	\begin{enumerate}[label=\alph*)]
		\item $(\overrightarrow{AB}, \overrightarrow{D'C'}) = 0^\circ$.
		\item $(\overrightarrow{BC}, \overrightarrow{D'A'}) = 0^\circ$.
		\item $(\overrightarrow{AC}, \overrightarrow{C'B'}) = 45^\circ$.
		\item $(\overrightarrow{AC}, \overrightarrow{A'B}) = 120^\circ$.
	\end{enumerate}
	
	\vspace{0.2cm}
	\noindent \textbf{Câu 13:} Trong không gian, cho hai vectơ $\vec{a}$ và $\vec{b}$ cùng có độ dài bằng $1$ và $(\vec{a}, \vec{b}) = 30^\circ$. Xét tính đúng sai của các mệnh đề sau:
	\begin{enumerate}[label=\alph*)]
		\item $\vec{a} \cdot \vec{b} = \frac{1}{2}$.
		\item $(\vec{a} + 2\vec{b})(\vec{a} - 3\vec{b}) = \frac{5\sqrt{3}}{2}$.
		\item $(\vec{a} - \vec{b})^2 = 0$.
		\item Hai vectơ $\vec{a} + \sqrt{3}\vec{b}$ và $\sqrt{3}\vec{a} - \vec{b}$ không vuông góc với nhau.
	\end{enumerate}
	
	\vspace{0.2cm}
	\noindent \textbf{Câu 14:} Cho tứ diện đều $ABCD$ cạnh bằng $1$. Các điểm $M, N$ lần lượt là trung điểm của $AB$ và $CD$. Xét tính đúng sai của các mệnh đề sau:
	\begin{enumerate}[label=\alph*)]
		\item $\overrightarrow{AB} \cdot \overrightarrow{AD} = 1$.
		\item $\overrightarrow{AB}(\overrightarrow{AB} + \overrightarrow{CA}) = \frac{1}{2}$.
		\item $AB \perp CD$.
		\item Góc giữa đường thẳng $MN$ với đường thẳng $BC$ bằng $30^\circ$.
	\end{enumerate}
	
	\vspace{0.2cm}
	\noindent \textbf{Câu 15:} Cho hình lập phương $ABCD.A'B'C'D'$. Xét tính đúng sai của các mệnh đề sau:
	\begin{enumerate}[label=\alph*)]
		\item $(\overrightarrow{AB}, \overrightarrow{C'D'}) = 180^\circ$.
		\item $(\overrightarrow{AD'}, \overrightarrow{B'C}) = 90^\circ$.
		\item $(\overrightarrow{A'C'}, \overrightarrow{BD}) = 90^\circ$.
		\item $(\overrightarrow{AC'}, \overrightarrow{A'D}) = 45^\circ$.
	\end{enumerate}
	
	\vspace{0.2cm}
	\noindent \textbf{Câu 16:} Trong không gian, cho hai vectơ $\vec{u}$ và $\vec{v}$ cùng có độ dài bằng $2$ và $(\vec{u}, \vec{v}) = 60^\circ$. Xét tính đúng sai của các mệnh đề sau:
	\begin{enumerate}[label=\alph*)]
		\item $\vec{u} \cdot \vec{v} = 2$.
		\item $|\vec{u} + \vec{v}| = 2\sqrt{3}$.
		\item $(\vec{u} - 2\vec{v})^2 = 12$.
		\item Hai vectơ $\vec{u} + 2\vec{v}$ và $2\vec{u} - \vec{v}$ vuông góc với nhau.
	\end{enumerate}
	
	\vspace{0.2cm}
	\noindent \textbf{Câu 17:} Cho tứ diện $OABC$ có $OA, OB, OC$ đôi một vuông góc và $OA = OB = OC = 1$. Gọi $M$ là trung điểm của $BC$. Xét tính đúng sai của các mệnh đề sau:
	\begin{enumerate}[label=\alph*)]
		\item $\overrightarrow{OB} \cdot \overrightarrow{OC} = 0$.
		\item $\overrightarrow{OM} = \frac{1}{2}(\overrightarrow{OB} + \overrightarrow{OC})$.
		\item $(\overrightarrow{OA}, \overrightarrow{OM}) = 45^\circ$.
		\item Độ dài của vectơ $\overrightarrow{OM}$ bằng $\frac{\sqrt{2}}{2}$.
	\end{enumerate}
	
	\vspace{0.2cm}
	\noindent \textbf{Câu 18:} Cho hình chóp $S.ABCD$ có đáy $ABCD$ là hình vuông cạnh $a$, $SA \perp (ABCD)$ và $SA = a$. Xét tính đúng sai của các mệnh đề sau:
	\begin{enumerate}[label=\alph*)]
		\item $\overrightarrow{AB} \cdot \overrightarrow{AD} = 0$.
		\item $\overrightarrow{SB} \cdot \overrightarrow{SD} = a^2$.
		\item $SC \perp BD$.
		\item Góc giữa hai vectơ $\overrightarrow{SA}$ và $\overrightarrow{SC}$ bằng $45^\circ$.
	\end{enumerate}
	
	\noindent \textbf{Câu 19:} Cho hình chóp $S.ABCD$ có tất cả các cạnh đều bằng $a$ và đáy $ABCD$ là hình vuông tâm $O$. Gọi $M$ là trung điểm của $CD$. Xét tính đúng sai của các mệnh đề sau:
	\begin{enumerate}[label=\alph*)]
		\item $\overrightarrow{OM} \cdot \overrightarrow{CB} = 0$.
		\item $\overrightarrow{SC} \cdot \overrightarrow{CB} = \overrightarrow{SC} \cdot \overrightarrow{CD}$.
		\item $\overrightarrow{MS} \cdot \overrightarrow{CB} = \frac{a^2}{2}$.
		\item $\cos(SD, BC) = \frac{1}{2}$.
	\end{enumerate}
	
	\vspace{0.2cm}
	\noindent \textbf{Câu 20:} Xét tính đúng, sai của các mệnh đề sau:
	\begin{enumerate}[label=\alph*)]
		\item $\overrightarrow{AB} + \overrightarrow{BC} = \overrightarrow{AC}$.
		\item $\overrightarrow{AB} + \overrightarrow{BC} + \overrightarrow{CD} = \overrightarrow{AD}$.
		\item $\overrightarrow{AB} + \overrightarrow{AD} = \overrightarrow{AC}$ với $ABCD$ là tứ giác.
		\item $\overrightarrow{AB} + \overrightarrow{AD} = \overrightarrow{AC}$ với $ABCD$ là hình bình hành.
	\end{enumerate}
	
	\vspace{0.2cm}
	\noindent \textbf{Câu 21:} Xét tính đúng, sai của các mệnh đề sau:
	\begin{enumerate}[label=\alph*)]
		\item $\overrightarrow{AB} - \overrightarrow{AC} = \overrightarrow{CB}$.
		\item $\overrightarrow{AB} - \overrightarrow{CB} = \overrightarrow{AC}$.
		\item $k\vec{a} = \vec{0} \Leftrightarrow \vec{a} = \vec{0}$.
		\item $\overrightarrow{AB} = k\overrightarrow{AC} \Leftrightarrow$ Ba điểm phân biệt $A, B, C$ thẳng hàng.
	\end{enumerate}
	
	\vspace{0.2cm}
	\noindent \textbf{Câu 22:} Xét tính đúng, sai của các mệnh đề sau:
	\begin{enumerate}[label=\alph*)]
		\item $k(\vec{a} + \vec{b}) = k\vec{a} + k\vec{b}$.
		\item Với $\vec{a} \neq \vec{0}$ thì $k\vec{a}$ cùng hướng với $\vec{a}$ nếu $k > 0$.
		\item $|k\vec{a}| = k|\vec{a}|$.
		\item Góc giữa hai vectơ $(\vec{u}, \vec{v}) \in [0^\circ; 180^\circ]$.
	\end{enumerate}
	
	\vspace{0.2cm}
	\noindent \textbf{Câu 23:} Xét tính đúng, sai của các mệnh đề sau:
	\begin{enumerate}[label=\alph*)]
		\item Với $\vec{a} \neq \vec{0}$ và $\vec{b} \neq \vec{0}$ thì $\vec{a} \cdot \vec{b} = |\vec{a}| \cdot |\vec{b}| \cdot \cos(\vec{a}, \vec{b})$.
		\item $(\vec{a})^2 = |\vec{a}|^2$.
		\item $\cos(\vec{a}, \vec{b}) = \frac{\vec{a} \cdot \vec{b}}{|\vec{a}| \cdot |\vec{b}|}$.
		\item Với $\vec{a} \neq \vec{0}$ và $\vec{b} \neq \vec{0}$ thì $\vec{a} \cdot \vec{b} = 0 \Leftrightarrow \vec{a} \perp \vec{b}$.
	\end{enumerate}
	
	\vspace{0.2cm}
	\noindent \textbf{Câu 24:} Xét tính đúng, sai của các mệnh đề sau:
	\begin{enumerate}[label=\alph*)]
		\item $\overrightarrow{MA} + \overrightarrow{MB} = 2\overrightarrow{MI}$ với $I$ là trung điểm đoạn $AB$ và điểm $M$ bất kì.
		\item $\overrightarrow{MA} + \overrightarrow{MB} + \overrightarrow{MC} = 3\overrightarrow{MG}$ với $G$ là trọng tâm $\triangle ABC$ và điểm $M$ bất kì.
		\item $\overrightarrow{MA} + \overrightarrow{MB} + \overrightarrow{MC} + \overrightarrow{MD} = 4\overrightarrow{MG}$ với $G$ là trọng tâm tứ diện $ABCD$ và điểm $M$ bất kì.
		\item Nếu $\overrightarrow{SB} + \overrightarrow{SD} = \overrightarrow{SA} + \overrightarrow{SC}$ thì chóp $S.ABCD$ có $ABCD$ là hình bình hành.
	\end{enumerate}
	
	\vspace{0.2cm}
	\noindent \textbf{Câu 25:} Cho hình hộp $ABCD.A'B'C'D'$. Xét tính đúng, sai của các mệnh đề sau:
	\begin{enumerate}[label=\alph*)]
		\item $\overrightarrow{AB} + \overrightarrow{CC'} = \overrightarrow{A'B'} + \overrightarrow{BB'}$.
		\item $\overrightarrow{AB} = \overrightarrow{CD}$.
		\item $\overrightarrow{AB} - \overrightarrow{BC'} = \overrightarrow{BD'}$.
		\item $\overrightarrow{AB} + \overrightarrow{AD} + \overrightarrow{AA'} = \overrightarrow{AC'}$.
	\end{enumerate}
	
	\vspace{0.2cm}
	\noindent \textbf{Câu 26:} Cho hình lăng trụ $ABC.A'B'C'$. Xét tính đúng, sai của các mệnh đề sau:
	\begin{enumerate}[label=\alph*)]
		\item $\overrightarrow{BA} + \overrightarrow{A'C'} = \overrightarrow{BC}$.
		\item Góc giữa $(\overrightarrow{BC}, \overrightarrow{AA'}) = (\overrightarrow{BC}, \overrightarrow{CC'}) = (\overrightarrow{BC}, \overrightarrow{BB'})$.
		\item $\overrightarrow{AB} + \overrightarrow{AA'} + \overrightarrow{B'C'} = \overrightarrow{AC'}$.
		\item $\overrightarrow{AA'} + \overrightarrow{B'C'} = \overrightarrow{AC'} - \overrightarrow{AB}$.
	\end{enumerate}
	
	\vspace{0.2cm}
	\noindent \textbf{Câu 27:} Cho tam giác $ABC$ có trọng tâm $G$ và $M$ là một điểm tùy ý. Xét tính đúng sai của các mệnh đề sau:
	\begin{enumerate}[label=\alph*)]
		\item $\overrightarrow{GA} + \overrightarrow{GB} + \overrightarrow{GC} = \vec{0}$.
		\item $\overrightarrow{MA} + \overrightarrow{MB} + \overrightarrow{MC} = 3\overrightarrow{MG}$.
		\item $\overrightarrow{AG} = \frac{1}{3}(\overrightarrow{AB} + \overrightarrow{AC})$.
		\item $\overrightarrow{GB} + \overrightarrow{GC} = \overrightarrow{GA}$.
	\end{enumerate}
	
	\vspace{0.2cm}
	\noindent \textbf{Câu 28:} Cho hình bình hành $ABCD$ tâm $O$ và điểm $S$ tùy ý. Xét tính đúng sai của các mệnh đề sau:
	\begin{enumerate}[label=\alph*)]
		\item $\overrightarrow{AB} + \overrightarrow{AD} = \overrightarrow{AC}$.
		\item $\overrightarrow{AB} - \overrightarrow{AD} = \overrightarrow{BD}$.
		\item $\overrightarrow{OA} + \overrightarrow{OB} + \overrightarrow{OC} + \overrightarrow{OD} = \vec{0}$.
		\item $\overrightarrow{SA} + \overrightarrow{SC} = \overrightarrow{SB} + \overrightarrow{SD}$.
	\end{enumerate}
	
	\vspace{0.2cm}
	\noindent \textbf{Câu 29:} Cho tứ diện $ABCD$. Gọi $M, N$ lần lượt là trung điểm của $AB$ và $CD$. Xét tính đúng sai của các mệnh đề sau:
	\begin{enumerate}[label=\alph*)]
		\item $\overrightarrow{MN} = \frac{1}{2}(\overrightarrow{AD} + \overrightarrow{BC})$.
		\item $\overrightarrow{MN} = \frac{1}{2}(\overrightarrow{AC} + \overrightarrow{BD})$.
		\item $\overrightarrow{AB} + \overrightarrow{BC} + \overrightarrow{CD} + \overrightarrow{DA} = \vec{0}$.
		\item $\overrightarrow{AC} + \overrightarrow{BD} = \overrightarrow{AD} + \overrightarrow{BC}$.
	\end{enumerate}
	
	\vspace{0.2cm}
	\noindent \textbf{Câu 30:} Cho hình lập phương $ABCD.A'B'C'D'$ cạnh $a$. Xét tính đúng sai của các mệnh đề sau:
	\begin{enumerate}[label=\alph*)]
		\item $\overrightarrow{AB} \cdot \overrightarrow{AD} = 0$.
		\item $\overrightarrow{AB} \cdot \overrightarrow{AA'} = a^2$.
		\item $\overrightarrow{AC'} = \overrightarrow{AB} + \overrightarrow{AD} + \overrightarrow{AA'}$.
		\item $(\overrightarrow{AD'}, \overrightarrow{B'C}) = 90^\circ$.
	\end{enumerate}
	
	\vspace{0.2cm}
	\noindent \textbf{Câu 31:} Trong không gian, cho hai vectơ $\vec{u}$ và $\vec{v}$ khác $\vec{0}$. Xét tính đúng sai của các mệnh đề sau:
	\begin{enumerate}[label=\alph*)]
		\item $\vec{u} \cdot \vec{v} > 0$ khi và chỉ khi góc $(\vec{u}, \vec{v})$ là góc nhọn.
		\item $\vec{u} \cdot \vec{v} < 0$ khi và chỉ khi góc $(\vec{u}, \vec{v})$ là góc tù.
		\item $\vec{u} \cdot \vec{v} = 0$ khi và chỉ khi hai vectơ vuông góc.
		\item $(\vec{u} + \vec{v})^2 = \vec{u}^2 + \vec{v}^2$ khi và chỉ khi $\vec{u} \perp \vec{v}$.
	\end{enumerate}
	
	\vspace{0.2cm}
	\noindent \textbf{Câu 32:} Cho hình chóp $S.ABCD$ có đáy $ABCD$ là hình chữ nhật và $SA \perp (ABCD)$. Xét tính đúng sai của các mệnh đề sau:
	\begin{enumerate}[label=\alph*)]
		\item $\overrightarrow{SA} \cdot \overrightarrow{AB} = 0$.
		\item $\overrightarrow{SA} \cdot \overrightarrow{CD} = 0$.
		\item $\overrightarrow{SD} \cdot \overrightarrow{BC} = 0$.
		\item $\overrightarrow{SC} \cdot \overrightarrow{BD} = 0$.
	\end{enumerate}
	
	\vspace{0.2cm}
	\noindent \textbf{Câu 33:} Cho hình lăng trụ tam giác $ABC.A'B'C'$. Xét tính đúng sai của các mệnh đề sau:
	\begin{enumerate}[label=\alph*)]
		\item $\overrightarrow{AB} + \overrightarrow{BC} + \overrightarrow{CC'} = \overrightarrow{AC'}$.
		\item $\overrightarrow{AA'} + \overrightarrow{BB'} + \overrightarrow{CC'} = 3\overrightarrow{AA'}$.
		\item $\overrightarrow{AB'} = \overrightarrow{AB} + \overrightarrow{AA'}$.
		\item $\overrightarrow{A'B} = \overrightarrow{AB} - \overrightarrow{AA'}$.
	\end{enumerate}
	
	\noindent \textbf{Câu 34:} Các mệnh đề sau đúng hay sai?
	\begin{enumerate}[label=\alph*)]
		\item Cho hình hộp $ABCDEFGH$. Tính tổng ba vectơ $\overrightarrow{AB} + \overrightarrow{AD} + \overrightarrow{AE} = \overrightarrow{AG}$.
		\item Cho hình hộp $ABCD.A'B'C'D'$. Các vectơ có điểm đầu và điểm cuối là các đỉnh của hình hộp và bằng vectơ $\overrightarrow{AB}$ là $\overrightarrow{DC}; \overrightarrow{A'B'}; \overrightarrow{D'C'}$.
		\item Cho hình hộp chữ nhật $ABCD.A'B'C'D'$. Khi đó, vectơ bằng vectơ $\overrightarrow{AB}$ là vectơ $\overrightarrow{D'C'}$.
		\item Cho tứ diện $ABCD$. Hỏi có bao nhiêu vectơ khác vectơ $\vec{0}$ mà mỗi vectơ có điểm đầu, điểm cuối là hai đỉnh của tứ diện $ABCD$ là 12.
	\end{enumerate}
	
	\vspace{0.2cm}
	\noindent \textbf{Câu 35:} Các mệnh đề sau đúng hay sai?
	\begin{enumerate}[label=\alph*)]
		\item Nếu trong ba vectơ $\vec{a}, \vec{b}, \vec{c}$ có một vectơ bằng $\vec{0}$ thì ba vectơ $\vec{a}, \vec{b}, \vec{c}$ đồng phẳng.
		\item Nếu giá của ba vectơ $\vec{a}, \vec{b}, \vec{c}$ cùng song song với một mặt phẳng thì ba vectơ $\vec{a}, \vec{b}, \vec{c}$ đồng phẳng.
		\item Nếu trong ba vectơ $\vec{a}, \vec{b}, \vec{c}$ có hai vectơ cùng phương thì ba vectơ $\vec{a}, \vec{b}, \vec{c}$ đồng phẳng.
		\item Nếu giá của ba vectơ $\vec{a}, \vec{b}, \vec{c}$ đồng quy thì ba vectơ $\vec{a}, \vec{b}, \vec{c}$ đồng phẳng.
	\end{enumerate}
	
	\vspace{0.2cm}
	\noindent \textbf{Câu 36:} Các mệnh đề sau đúng hay sai?
	\begin{enumerate}[label=\alph*)]
		\item Cho $\vec{a}, \vec{b}, \vec{c}$ đều khác $\vec{0}$. Ba vectơ $\vec{a}, \vec{b}, \vec{c}$ đồng phẳng khi và chỉ khi giá của chúng cùng nằm trên một mặt phẳng.
		\item Với tứ diện $ABCD$ bất kì ta luôn có $\overrightarrow{AC} + \overrightarrow{BD} = \overrightarrow{AD} + \overrightarrow{BC}$.
		\item Một đường thẳng cắt hai đường thẳng cho trước thì tồn tại một mặt phẳng chứa cả ba đường thẳng đó.
		\item Với hình hộp $ABCD.A'B'C'D'$ bất kì ta luôn có $\overrightarrow{AB} + \overrightarrow{AD} + \overrightarrow{AA'} = \overrightarrow{C'A}$.
	\end{enumerate}
	
	
	% --- BÀI TẬP TRẢ LỜI NGẮN ---
	\section{Câu hỏi trả lời ngắn}
	
	\noindent \textbf{Câu 1:} Cho hình lăng trụ tam giác $ABC.A'B'C'$ biết tam giác $ABC$ đều có cạnh bằng $3$. Tính độ dài của vectơ $\overrightarrow{CA} + \overrightarrow{B'C'}$.
	
	\vspace{0.2cm}
	\noindent \textbf{Câu 2:} Cho hình lập phương $ABCD.A'B'C'D'$. Tính số đo góc giữa hai vectơ $\overrightarrow{AB}$ và $\overrightarrow{A'C'}$ (theo đơn vị độ).
	
	\vspace{0.2cm}
	\noindent \textbf{Câu 3:} Cho tứ diện $S.ABC$ có các cạnh $SA = SB = SC = AB = AC = 2$ và $BC = 2\sqrt{2}$. Tính tích vô hướng của hai vectơ $\overrightarrow{SC} \cdot \overrightarrow{AB}$.
	
	\vspace{0.2cm}
	\noindent \textbf{Câu 4:} Cho tứ diện $ABCD$. Gọi $E, F$ lần lượt là trung điểm của $AB$ và $CD$. Cho $AB = 2a, CD = 2b, EF = 2c$. Với $M$ là một điểm tùy ý trong không gian, biết hệ thức liên hệ độ dài luôn thỏa mãn $MA^2 + MB^2 = k \cdot ME^2 + l \cdot a^2$. Tính giá trị của biểu thức $k + l$.
	
	\noindent \textbf{Câu 5:} Cho tứ diện $ABCD$ có các điểm $M, N, P$ lần lượt thuộc các cạnh $BC, BD$ và $AC$ sao cho $BC = 4BM$, $AC = 3AP$, $BD = 2BN$. Mặt phẳng $(MNP)$ cắt đường thẳng $AD$ tại điểm $Q$. Tính tỉ số $\frac{AQ}{AD}$.
	
	\vspace{0.2cm}
	\noindent \textbf{Câu 6:} Cho tứ diện $ABCD$. Lấy các điểm $M, N$ lần lượt thuộc các cạnh $AB, AC$ sao cho $AB = 3AM$, $AC = 4AN$. Gọi $P$ là trung điểm của cạnh $CD$. Mặt phẳng $(MNP)$ cắt đường thẳng $BD$ tại điểm $Q$. Tính tỉ số $\frac{BQ}{BD}$.
	
	\noindent \textbf{Câu 7:} Cho tứ diện đều $ABCD$ có cạnh bằng $a$ và $M$ là trung điểm của $CD$. Tính góc $(\overrightarrow{AB}, \overrightarrow{CD})$.
	\begin{figure}[H]
		\centering
		\begin{tikzpicture}[scale=1.2, line join=round, line cap=round]
			\coordinate (B) at (0,0);
			\coordinate (C) at (3,-0.5);
			\coordinate (D) at (4,1.2);
			\coordinate (A) at (2,3);
			% Tính tọa độ trung điểm M của CD thủ công để tránh lỗi thư viện calc
			\coordinate (M) at (3.5, 0.35); 
			
			\draw[dashed] (B) -- (D);
			\draw[dashed] (B) -- (M);
			
			\draw (A) -- (B) -- (C) -- (D) -- cycle;
			\draw (A) -- (C);
			\draw (A) -- (M);
			
			\draw[-stealth, thick, blue] (A) -- (B);
			\draw[-stealth, thick, blue] (C) -- (D);
			
			\node[above] at (A) {$A$};
			\node[below left] at (B) {$B$};
			\node[below] at (C) {$C$};
			\node[right] at (D) {$D$};
			\node[right] at (M) {$M$};
		\end{tikzpicture}
	\end{figure}
	
	\vspace{0.2cm}
	\noindent \textbf{Câu 8:} Cho hình lập phương $ABCD.A'B'C'D'$. Tính góc giữa hai vectơ $\overrightarrow{BD}$ và $\overrightarrow{B'C}$.
	\begin{figure}[H]
		\centering
		\begin{tikzpicture}[scale=1.2, line join=round, line cap=round]
			\coordinate (A) at (1.5,1.2);
			\coordinate (B) at (0,0);
			\coordinate (C) at (3,-0.2);
			\coordinate (D) at (4.5,1);
			\coordinate (A') at (1.5,4.2);
			\coordinate (B') at (0,3);
			\coordinate (C') at (3,2.8);
			\coordinate (D') at (4.5,4);
			
			\draw[dashed] (A) -- (B);
			\draw[dashed] (A) -- (D);
			\draw[dashed] (A) -- (A');
			
			\draw[dashed] (B') -- (D);
			\draw[dashed, thick, blue, -stealth] (B) -- (D);
			
			\draw (B) -- (C) -- (D) -- (D') -- (C') -- (B') -- cycle;
			\draw (C) -- (C');
			\draw (B') -- (A') -- (D');
			
			\draw[-stealth, thick, blue] (B') -- (C);
			
			\node[left] at (A) {$A$};
			\node[below left] at (B) {$B$};
			\node[below right] at (C) {$C$};
			\node[right] at (D) {$D$};
			\node[left] at (A') {$A'$};
			\node[left] at (B') {$B'$};
			\node[right] at (C') {$C'$};
			\node[right] at (D') {$D'$};
		\end{tikzpicture}
	\end{figure}
	
	\vspace{0.2cm}
	\noindent \textbf{Câu 9:} Cho tứ diện $OABC$ có các cạnh $OA, OB, OC$ đôi một vuông góc và $OA = OB = OC = 1$. Gọi $M$ là trung điểm của cạnh $AB$. Tính góc giữa hai vectơ $\overrightarrow{OM}$ và $\overrightarrow{AC}$.
	\begin{figure}[H]
		\centering
		\begin{tikzpicture}[scale=1.2, line join=round, line cap=round]
			\coordinate (O) at (0,0);
			\coordinate (C) at (0,3);
			\coordinate (A) at (2,-1.5);
			\coordinate (B) at (3.5,0.5);
			% Tính tọa độ trung điểm M của AB thủ công
			\coordinate (M) at (2.75, -0.5); 
			
			\draw[dashed] (O) -- (B);
			\draw[dashed, thick, blue, -stealth] (O) -- (M);
			
			\draw (O) -- (C);
			\draw (O) -- (A);
			\draw (A) -- (B);
			\draw (C) -- (A);
			\draw (C) -- (B);
			
			\draw[-stealth, thick, blue] (A) -- (C);
			
			\node[left] at (O) {$O$};
			\node[above] at (C) {$C$};
			\node[below] at (A) {$A$};
			\node[right] at (B) {$B$};
			\node[below right] at (M) {$M$};
		\end{tikzpicture}
	\end{figure}
	
	\vspace{0.2cm}
	\noindent \textbf{Câu 10:} Cho hình chóp $S.ABCD$ có đáy $ABCD$ là hình bình hành và mặt bên $SAB$ là tam giác đều. Tính góc giữa hai vectơ $\overrightarrow{DC}$ và $\overrightarrow{BS}$.
	\begin{figure}[H]
		\centering
		\begin{tikzpicture}[scale=1.2, line join=round, line cap=round]
			\coordinate (B) at (0,0);
			\coordinate (C) at (3,0);
			\coordinate (A) at (1.5,1);
			\coordinate (D) at (4.5,1);
			\coordinate (S) at (0.5,2.5);
			\coordinate (E) at (-1.5,-1);
			
			\draw[dashed] (B) -- (A) -- (D);
			\draw[dashed] (S) -- (A);
			
			\draw (S) -- (B) -- (C) -- (D) -- cycle;
			\draw (S) -- (C);
			
			\draw[thick, gray, -stealth] (B) -- (E) node[below left, text=black] {$E$};
			
			\draw[-stealth, thick, blue] (D) -- (C);
			\draw[-stealth, thick, blue] (B) -- (S);
			
			% Cung góc 60 độ (giữa BA và BS)
			\draw (B) ++(33.7:0.5) arc (33.7:78.7:0.5);
			\node at (0.8, 0.6) {$60^\circ$};
			
			% Cung góc 120 độ (giữa BE và BS)
			\draw (B) ++(78.7:0.6) arc (78.7:213.7:0.6);
			\node at (-0.5, 0.8) {$120^\circ$};
			
			\node[above] at (S) {$S$};
			\node[below left] at (B) {$B$};
			\node[below right] at (C) {$C$};
			\node[right] at (D) {$D$};
			\node[above right] at (A) {$A$};
		\end{tikzpicture}
	\end{figure}
	
	\noindent \textbf{Câu 11:} Cho hình chóp $S.ABCD$ có đáy $ABCD$ là hình bình hành. Mặt bên $ASB$ là tam giác vuông cân tại $S$ và có cạnh $AB = a$. Gọi $M$ là trung điểm của $AB$. Hãy tính: $\overrightarrow{DC} \cdot \overrightarrow{AS}$.
	\begin{figure}[H]
		\centering
		\begin{tikzpicture}[scale=1.2, line join=round, line cap=round]
			\coordinate (B) at (0,0);
			\coordinate (C) at (3,-0.2);
			\coordinate (A) at (1.5,1.2);
			\coordinate (D) at (4.5,1);
			\coordinate (S) at (1.5,3.5);
			\coordinate (M) at (0.75,0.6); % Trung điểm AB
			\coordinate (E) at (-1.5,-1.2); % Tia BE
			
			% Các nét khuất
			\draw[dashed] (B) -- (A) -- (D);
			\draw[dashed] (S) -- (A);
			\draw[dashed] (S) -- (M);
			
			% Các nét thấy
			\draw (S) -- (B) -- (C) -- (D) -- cycle;
			\draw (S) -- (C);
			
			% Tia BE
			\draw[thick, gray, -stealth] (B) -- (E) node[below left, text=black] {$E$};
			
			% Các vectơ màu xanh
			\draw[-stealth, thick, blue] (D) -- (C);
			\draw[-stealth, thick, blue, dashed] (A) -- (S);
			
			% Cung chỉ góc
			\draw (B) ++(38.6:0.5) arc (38.6:66.8:0.5);
			\node at (0.8, 0.7) {\footnotesize $45^\circ$};
			
			\draw (B) ++(66.8:0.6) arc (66.8:218.6:0.6);
			\node at (-0.5, 0.5) {\footnotesize $135^\circ$};
			
			% Ký hiệu vuông góc tại S
			\draw[dashed] (1.3, 3.1) -- (1.5, 2.95) -- (1.65, 3.15);
			
			% Đặt tên các đỉnh
			\node[above] at (S) {$S$};
			\node[below right] at (B) {$B$};
			\node[below right] at (C) {$C$};
			\node[right] at (D) {$D$};
			\node[above right] at (A) {$A$};
			\node[below right] at (M) {$M$};
		\end{tikzpicture}
	\end{figure}
	
	\vspace{0.2cm}
	\noindent \textbf{Câu 12:} Cho hình chóp $S.ABCD$ có đáy $ABCD$ là hình bình hành. Mặt bên $ASB$ là tam giác vuông cân tại $S$ và có cạnh $AB = a$. Gọi $M$ là trung điểm của $AB$. Hãy tính: $\overrightarrow{DC} \cdot \overrightarrow{MS}$.
	\begin{figure}[H]
		\centering
		\begin{tikzpicture}[scale=1.2, line join=round, line cap=round]
			\coordinate (B) at (0,0);
			\coordinate (C) at (3,-0.2);
			\coordinate (A) at (1.5,1.2);
			\coordinate (D) at (4.5,1);
			\coordinate (S) at (1.5,3.5);
			\coordinate (M) at (0.75,0.6); 
			\coordinate (E) at (-1.5,-1.2); 
			
			\draw[dashed] (B) -- (A) -- (D);
			\draw[dashed] (S) -- (A);
			
			\draw (S) -- (B) -- (C) -- (D) -- cycle;
			\draw (S) -- (C);
			
			\draw[thick, gray, -stealth] (B) -- (E) node[below left, text=black] {$E$};
			
			\draw[-stealth, thick, blue] (D) -- (C);
			% Vẽ vectơ MS đè lên nét đứt
			\draw[-stealth, thick, blue, dashed] (M) -- (S);
			
			\draw (B) ++(38.6:0.5) arc (38.6:66.8:0.5);
			\node at (0.8, 0.7) {\footnotesize $45^\circ$};
			
			\draw (B) ++(66.8:0.6) arc (66.8:218.6:0.6);
			\node at (-0.5, 0.5) {\footnotesize $135^\circ$};
			
			\draw[dashed] (1.3, 3.1) -- (1.5, 2.95) -- (1.65, 3.15);
			
			\node[above] at (S) {$S$};
			\node[below right] at (B) {$B$};
			\node[below right] at (C) {$C$};
			\node[right] at (D) {$D$};
			\node[above right] at (A) {$A$};
			\node[below right] at (M) {$M$};
		\end{tikzpicture}
	\end{figure}
	
	\vspace{0.2cm}
	\noindent \textbf{Câu 13:} Cho hình chóp $S.ABC$ có $SA = SB = SC = AB = AC = a$ và $BC = a\sqrt{2}$. Tính góc giữa các vectơ $\overrightarrow{SC}$ và $\overrightarrow{AB}$.
	\begin{figure}[H]
		\centering
		\begin{tikzpicture}[scale=1.2, line join=round, line cap=round]
			\coordinate (A) at (1,1);
			\coordinate (B) at (2.5,-1.5);
			\coordinate (C) at (5,1.5);
			\coordinate (S) at (2.5,3.5);
			\coordinate (E) at (-0.5,0.6); % Tia dọc theo tia đối CA
			
			\draw[dashed] (A) -- (C);
			\draw[thick, gray, -stealth] (A) -- (E);
			
			\draw (S) -- (A) -- (B) -- (C) -- cycle;
			\draw (S) -- (B);
			
			% Ký hiệu vuông góc tại A
			\draw[dashed] (1.2, 0.65) -- (1.5, 0.75) -- (1.3, 1.1);
			
			% Góc 120 độ
			\draw (A) ++(194.9:0.4) arc (194.9:300.9:0.4);
			\node at (0.5, 0.1) {\footnotesize $120^\circ$};
			
			% Tên đỉnh
			\node[above] at (S) {$S$};
			\node[above right] at (A) {$A$};
			\node[below] at (B) {$B$};
			\node[right] at (C) {$C$};
			
			% Các cạnh có giá trị
			\node[left] at (1.75, 2.25) {$a$};
			\node[right] at (2.5, 1) {$a$};
			\node[right] at (3.75, 2.5) {$a$};
			\node[below left] at (1.75, -0.25) {$a$};
			\node[above] at (3, 1.25) {$a$};
			\node[below right] at (3.75, 0) {$a\sqrt{2}$};
		\end{tikzpicture}
	\end{figure}
	
	\newpage
	
	% =====================
	% PHẦN 5: TỔNG KẾT
	% =====================
	\section{Tổng kết bài}
	
	\begin{tcolorbox}[colback=yellow!10!white, colframe=orange!60!black,
		fonttitle=\bfseries, title={Sơ đồ hệ thống hóa \& Công thức cần nhớ}]
		\begin{itemize}
			\item \textbf{Quy tắc 3 điểm:} $\vec{AB} + \vec{BC} = \vec{AC}$
			\item \textbf{Quy tắc hình bình hành:} Nếu $ABCD$ là hình bình hành thì $\vec{AB} + \vec{AD} = \vec{AC}$
			\item \textbf{Quy tắc hình hộp:} Nếu $ABCD.A'B'C'D'$ là hình hộp thì $\vec{AB} + \vec{AD} + \vec{AA'} = \vec{AC'}$
			\item \textbf{Tích vô hướng:} $\vec{a} \cdot \vec{b} = |\vec{a}| \cdot |\vec{b}| \cdot \cos(\vec{a}, \vec{b})$
			\item \textbf{Sự đồng phẳng:} $\vec{c} = m\vec{a} + n\vec{b}$ (với $\vec{a}, \vec{b}$ không cùng phương).
		\end{itemize}
	\end{tcolorbox}
	
	\vspace{5mm}
	\textbf{Bài tập về nhà:}
	\begin{itemize}
		\item Hoàn thành các bài tập phần Tự luận (thuộc khung màu tương ứng với năng lực) chưa giải quyết trên lớp.
		\item Làm toàn bộ các câu Trắc nghiệm, Đúng/Sai và Trả lời ngắn còn lại.
	\end{itemize}
	
\end{document}